\documentclass[conference]{IEEEtran}
\IEEEoverridecommandlockouts
% The preceding line is only needed to identify funding in the first footnote. If that is unneeded, please comment it out.
\usepackage{cite}
\usepackage{amsmath,amssymb,amsfonts}
% \usepackage{algorithmic}
\usepackage{graphicx}
\usepackage{textcomp}
\usepackage{xcolor}
% \usepackage[margin=1in]{geometry}
\usepackage{amsmath, amssymb, amsthm}
\usepackage{graphicx}
\usepackage{booktabs}
% \usepackage{caption}
\usepackage{subcaption}
\usepackage{algorithm}
% \usepackage{multirow}
\usepackage{algpseudocode}
\usepackage{hyperref}
\usepackage{pgfplots}
\usepackage{pgfplotstable}

\pgfplotsset{compat=1.18}
\def\BibTeX{{\rm B\kern-.05em{\sc i\kern-.025em b}\kern-.08em
    T\kern-.1667em\lower.7ex\hbox{E}\kern-.125emX}}
\begin{document}

\title{TRINITY: A Tri-order Regularized Optimization with Adaptivity}

% \author{\IEEEauthorblockN{1\textsuperscript{st} Given Name Surname}
% \IEEEauthorblockA{\textit{dept. name of organization (of Aff.)} \\
% \textit{name of organization (of Aff.)}\\
% City, Country \\
% email address or ORCID}
% \and
% \IEEEauthorblockN{2\textsuperscript{nd} Given Name Surname}
% \IEEEauthorblockA{\textit{dept. name of organization (of Aff.)} \\
% \textit{name of organization (of Aff.)}\\
% City, Country \\
% email address or ORCID}
% \and
% \IEEEauthorblockN{3\textsuperscript{rd} Given Name Surname}
% \IEEEauthorblockA{\textit{dept. name of organization (of Aff.)} \\
% \textit{name of organization (of Aff.)}\\
% City, Country \\
% email address or ORCID}
% \and
% \IEEEauthorblockN{4\textsuperscript{th} Given Name Surname}
% \IEEEauthorblockA{\textit{dept. name of organization (of Aff.)} \\
% \textit{name of organization (of Aff.)}\\
% City, Country \\
% email address or ORCID}
% \and
% \IEEEauthorblockN{5\textsuperscript{th} Given Name Surname}
% \IEEEauthorblockA{\textit{dept. name of organization (of Aff.)} \\
% \textit{name of organization (of Aff.)}\\
% City, Country \\
% email address or ORCID}
% \and
% \IEEEauthorblockN{6\textsuperscript{th} Given Name Surname}
% \IEEEauthorblockA{\textit{dept. name of organization (of Aff.)} \\
% \textit{name of organization (of Aff.)}\\
% City, Country \\
% email address or ORCID}
% }

\maketitle

\begin{abstract}
The current state of the art in deep learning optimization typically employs zeroth-, first-, and second-order gradient updates in isolation. We propose Trinity, a novel optimization framework that unifies these approaches by using zeroth-order curvature sensing for selective natural-gradient preconditioning. Trinity introduces a forward-only second-difference curvature sensor that estimates the effective rank of the Fisher Information Matrix per block at minimal overhead ($\sim$6\%). This sensor output drives a controller that restricts expensive K-FAC inversions to layers where the geometry is highly curved, while an AdamW actuator with magnitude grafting ensures stable scale-safe updates elsewhere. Additionally, the zeroth-order signal acts as a direct saddle-escape actuator. Benchmark evaluations on extreme 100:1 long-tailed label imbalances across diverse architectures demonstrate that Trinity matches the performance of always-on K-FAC at a fraction of the inversion cost.
\end{abstract}

\begin{IEEEkeywords}
Optimizer, K-FAC, Minimization, CNN, Vision Transformer, Gradient Centralization, Preconditioning
\end{IEEEkeywords}



\section{Introduction}
\label{sec:introduction}

First-order optimization methods, though empirically scalable, are fundamentally geometry-blind. Second-order methods, such as natural gradient descent, address this by preconditioning updates with curvature information, but they pay an $\mathcal{O}(n^3)$ cost per block per inversion interval. The standard response to this computational bottleneck is a global schedule, where all layers are inverted uniformly every $\tau$ steps.

However, the value of preconditioning is not uniform across blocks. If a block's Fisher Information Matrix is near-isotropic, the preconditioned gradient $\mathbf{F}^{-1}\nabla L$ is roughly proportional to the raw gradient $\nabla L$, meaning the expensive matrix inversion buys almost nothing. Despite this, such anisotropy is rarely measured during training, and the inversion budget is spent uniformly across the network.

The primary reason this anisotropy remains unmeasured is that traditional diagnostics require Hessian-vector products or full eigendecompositions. These operations are typically as expensive as the preconditioning step they are meant to evaluate, making them impractical for dynamic runtime control.

Our core insight is that two forward passes can provide a second difference rather than a first difference, yielding an unbiased quadratic form $\mathbf{z}^T \mathbf{H} \mathbf{z}$. This operates at the same computational cost as a standard zeroth-order gradient estimate but is significantly more informative, circumventing the $\mathcal{O}(d)$ variance issues that dominate standard backpropagation in differentiable settings.

Building on this, we propose \textbf{Trinity}: a unified framework where zeroth-order (ZO) signals sense the curvature, while second-order (K-FAC) and first-order (AdamW) mechanisms act on it. Furthermore, the ZO component acts directly as a saddle-escape actuator. All three orders are present, each functioning in a structurally defensible role.

Our concrete contributions are:
\begin{enumerate}
    \item \textbf{A forward-only curvature sensor:} Using RDSA second differences, we obtain unbiased estimates of $\mathrm{tr}(\mathbf{H})$ and $\|\mathbf{H}\|_F^2$ per parameter block, yielding a normalized effective rank $\rho \in (0,1]$ at $\sim$6\% overhead.
    \item \textbf{A measurement study of Fisher anisotropy:} We empirically characterize where the metric is actually non-Euclidean across layers and training time under severe label imbalance.
    \item \textbf{Selective natural-gradient preconditioning (Trinity):} A per-block controller spends a bounded K-FAC inversion budget only where the sensor indicates curved geometry, falling back to AdamW elsewhere with scale-safe magnitude grafting.
    \item \textbf{Cost result:} Trinity matches the performance of always-on K-FAC at matched wall-clock time while inverting only a bounded fraction of blocks.
\end{enumerate}

\section{Related Work}

\subsection{First-Order Adaptive Methods and Structural Regularization}
Stochastic Gradient Descent and its momentum-based variants have historically underpinned deep neural network training. To address sensitivity to hyperparameter tuning and ill-conditioned landscapes, adaptive first-order methods like Adam \cite{kingma2014adam} and RMSProp were introduced. The generalization gap of Adam was fundamentally addressed by Loshchilov and Hutter \cite{loshchilov2019decoupled} with AdamW, which decoupled weight decay from the adaptive gradient updates. Furthering first-order stability, structural constraints applied directly to the gradient tensor, such as Gradient Centralization (GC) \cite{yong2020gradient}, project weight gradients to have a zero mean across non-output dimensions. While these first-order enhancements drastically improve empirical stability, they remain fundamentally oblivious to the underlying curvature of the loss surface.

\subsection{Second-Order Preconditioning and Grafting}
To overcome geometric blindness, second-order optimization explicitly accounts for landscape curvature. Amari's Natural Gradient Descent (NGD) \cite{amari1998natural} adjusts updates using the inverse Fisher Information Matrix (FIM). To achieve tractable scaling, Martens and Grosse \cite{martens2015optimizing} proposed Kronecker-Factored Approximate Curvature (K-FAC) for linear layers, later extended to convolutional layers (KFC) by Grosse and Martens. Extensions of K-FAC have been successfully applied to massive architectures \cite{osawa2019large}. Other recent preconditioning approaches include Shampoo (Gupta et al. 2018), AdaHessian, SOAP (Vyas et al. 2024), and Muon. To safely combine the directional benefits of second-order steps with the robust step sizes of first-order methods, grafting (Agarwal et al.) has emerged as a crucial technique.

\subsection{Zeroth-Order Optimization and Perturbation}
Zeroth-Order (ZO) optimization bypasses analytical backpropagation, constructing trajectories strictly from forward objective evaluations. Methods range from Spall's SPSA \cite{spall1992multivariate} to Nesterov and Spokoiny's Gaussian smoothing estimators \cite{nesterov2017random}. Recently, ZO principles have been integrated into deep learning as landscape-aware regularizers or memory-efficient optimizers, such as MeZO (Malladi et al. 2023). Crucially, Sharpness-Aware Minimization (SAM) (Foret et al. 2021) established the motivation of escaping sharp minima via an extra forward/backward pass. Additionally, perturbed gradient descent (Ge et al. 2015, Jin et al. 2017) demonstrated that injecting random noise allows the optimizer to efficiently escape saddle points.

\subsection{Adaptive Curvature Use}
The classical ancestor of curvature-driven adaptation is the Levenberg-Marquardt (LM) gain ratio, which adapts damping based on the agreement between actual and predicted loss reduction. K-FAC employs an LM-style damping adaptation. We position our approach explicitly against such global, computationally heavy evaluations: our proposed ZO sensing signal is forward-only, dimension-normalized, and evaluated per-block, directly measuring structural anisotropy without requiring expensive matrix inversions or Hessian-vector products.

\section{Background and Notation}
\label{sec:background}

In this section, we establish the mathematical notation, foundational optimization paradigms, and regularized objective functions that form the basis of our work.

\subsection{Problem Setup and Core Notation}
\label{subsec:notation}

We consider a deep neural network parameterized by a weight vector $\boldsymbol{\theta} \in \mathbb{R}^{d}$. The network defines a non-linear mapping from an input token or image $\mathbf{x} \in \mathcal{X}$ to a vector of class logits $\mathbf{z}(\mathbf{x}; \boldsymbol{\theta}) \in \mathbb{R}^{C}$, where $C$ denotes the number of target categories. The conditional probability distribution over the classes is modeled via the softmax operator:
\begin{equation}
p_k(\mathbf{x}; \boldsymbol{\theta}) = \frac{\exp\left(z_k(\mathbf{x}; \boldsymbol{\theta})\right)}{\sum_{j=1}^{C} \exp\left(z_j(\mathbf{x}; \boldsymbol{\theta})\right)}, \qquad \forall k \in \{1, \dots, C\}.
\label{eq:softmax}
\end{equation}

Given a stochastic mini-batch $\mathcal{B} = \{(\mathbf{x}_i, y_i)\}_{i=1}^{B}$ drawn from the training distribution, where $y_i \in \{1, \dots, C\}$ represents the ground-truth label, the empirical risk minimizer seeks to minimize the batch loss function:
\begin{equation}
\mathcal{L}(\boldsymbol{\theta}; \mathcal{B}) = \frac{1}{B}\sum_{i=1}^{B}\ell(\mathbf{x}_i, y_i; \boldsymbol{\theta}),
\label{eq:batch_loss}
\end{equation}
where $\ell(\cdot)$ is the per-sample objective. At optimization step $t$, the standard first-order backpropagated gradient evaluated over the active mini-batch $\mathcal{B}^{(t)}$ is defined as:
\begin{equation}
\mathbf{g}^{(t)} = \nabla_{\boldsymbol{\theta}} \mathcal{L}(\boldsymbol{\theta}^{(t)}; \mathcal{B}^{(t)}).
\label{eq:full_grad}
\end{equation}

For a layer-wise analysis, we define a block partition of the parameters $\boldsymbol{\theta} = \{\boldsymbol{\theta}^{(\ell)}\}$, with corresponding sizes $d_\ell$. We denote the weight matrix and bias vector of layer $\ell$ by $\mathbf{W}^{(\ell)}$ and $\mathbf{b}^{(\ell)}$, respectively. A block is considered eligible for second-order preconditioning if its parameter tensor has dimension $\ge 2$ (e.g., weights of Linear or Convolutional layers). One-dimensional blocks, such as normalization gains or standalone biases, are permanently restricted to first-order updates. We denote the incoming layer activations by $\mathbf{A}^{(\ell)}$ and the backpropagated pre-activation sensitivities by $\mathbf{S}^{(\ell)}$.

\subsection{First-Order Optimization and Gradient Operations}
\label{subsec:first_order_bg}

\paragraph{Adaptive Moment Estimation (Adam/AdamW)} First-order adaptive optimizers track temporal variations in gradients using exponential moving averages. Standard Adam maintains a first moment vector $\mathbf{m}^{(t)}$ (momentum) and an uncentered second moment vector $\mathbf{v}^{(t)}$ (variance):
\begin{align}
\mathbf{m}^{(t)} &= \beta_1 \mathbf{m}^{(t-1)} + (1-\beta_1)\mathbf{g}^{(t)}, \\
\mathbf{v}^{(t)} &= \beta_2 \mathbf{v}^{(t-1)} + (1-\beta_2)\bigl(\mathbf{g}^{(t)} \odot \mathbf{g}^{(t)}\bigr),
\label{eq:moments}
\end{align}
where $\beta_1, \beta_2 \in (0, 1)$ are decay hyperparameters, and $\odot$ denotes the Hadamard product. Applying bias-correction yields $\hat{\mathbf{m}}^{(t)} = \mathbf{m}^{(t)} / (1-\beta_1^t)$ and $\hat{\mathbf{v}}^{(t)} = \mathbf{v}^{(t)} / (1-\beta_2^t)$. The decoupled weight decay variant, AdamW, separates $L_2$ regularization from these adaptive moment trajectories, yielding the update rule:
\begin{equation}
\boldsymbol{\theta}^{(t+1)} = \boldsymbol{\theta}^{(t)} - \eta^{(t)} \left( \frac{\hat{\mathbf{m}}^{(t)}}{\sqrt{\hat{\mathbf{v}}^{(t)}}+\epsilon} + \lambda_{\mathrm{wd}}\boldsymbol{\theta}^{(t)} \right),
\label{eq:adamw_update}
\end{equation}
where $\eta^{(t)}$ is the learning rate, $\epsilon > 0$ is a stabilization constant, and $\lambda_{\mathrm{wd}}$ is the weight decay coefficient.

\paragraph{Gradient Centralization (GC)} Gradient Centralization acts as a structural constraint directly on the backpropagated gradients of multi-dimensional tensors (e.g., weights of Linear or Conv2d layers). For a weight tensor $\mathbf{W} \in \mathbb{R}^{n \times d_1 \times \cdots \times d_r}$ with computed gradient $\mathbf{G}$, GC projects the gradient matrix by subtracting its mean along the non-output dimensions:
\begin{equation}
\tilde{G}_{i, j_1, \dots, j_r} = G_{i, j_1, \dots, j_r} - \frac{1}{\prod_{k=1}^r d_k} \sum_{k_1,\dots,k_r} G_{i, k_1, \dots, k_r}.
\label{eq:gradient_centralization}
\end{equation}

\subsection{Curvature-Aware Optimization via K-FAC}
\label{subsec:kfac_bg}

Natural Gradient Descent (NGD) scales the gradient by the inverse of the Fisher Information Matrix (FIM), denoted as $\mathbf{F}(\boldsymbol{\theta})$:
\begin{equation}
\Delta \boldsymbol{\theta}_{\mathrm{NGD}} = -\eta \mathbf{F}(\boldsymbol{\theta})^{-1} \nabla_{\boldsymbol{\theta}} \mathcal{L}(\boldsymbol{\theta}).
\label{eq:ngd}
\end{equation}
As explicitly calculating and inverting the full $d \times d$ matrix $\mathbf{F}$ is computationally expensive, Kronecker-Factored Approximate Curvature (K-FAC) introduces a block-diagonal approximation. For a structural layer $\ell$, the corresponding diagonal block $\mathbf{F}^{(\ell)}$ is factored into the Kronecker product ($\otimes$) of two significantly smaller covariance matrices:
\begin{equation}
\begin{aligned}
\mathbf{F}^{(\ell)} &\approx \mathbf{A}^{(\ell)} \otimes \mathbf{S}^{(\ell)}, \\
\text{where } \mathbf{A}^{(\ell)} &= \mathbb{E}\bigl[\mathbf{A}^{(\ell)\top}\mathbf{A}^{(\ell)}\bigr], \\
\mathbf{S}^{(\ell)} &= \mathbb{E}\bigl[\mathbf{S}^{(\ell)\top}\mathbf{S}^{(\ell)}\bigr].
\end{aligned}
\end{equation}
Using empirical mini-batches, the running factors $\hat{\mathbf{A}}^{(\ell)}_t$ and $\hat{\mathbf{S}}^{(\ell)}_t$ are tracked via an exponential moving average with decay $\rho \in (0, 1)$:
\begin{equation}
\label{eq:kfac_ema}
\begin{aligned}
\hat{\mathbf{A}}^{(\ell)}_t &= \rho \hat{\mathbf{A}}^{(\ell)}_{t-1} + (1-\rho)\mathbf{A}^{(\ell)}_{\text{emp},t}, \\
\hat{\mathbf{S}}^{(\ell)}_t &= \rho \hat{\mathbf{S}}^{(\ell)}_{t-1} + (1-\rho)\mathbf{S}^{(\ell)}_{\text{emp},t}.
\end{aligned}
\end{equation}

To maintain numerical stability during matrix inversion, a trace-aware balancing scalar $\pi^{(\ell)}$ scales the Tikhonov damping coefficient $\lambda^{(t)}$:
\begin{equation}
\pi^{(\ell)} = \sqrt{\frac{\mathrm{tr}(\hat{\mathbf{A}}^{(\ell)}_t)\,\dim(\hat{\mathbf{S}}^{(\ell)}_t)}{\mathrm{tr}(\hat{\mathbf{S}}^{(\ell)}_t)\,\dim(\hat{\mathbf{A}}^{(\ell)}_t) + \delta}},
\label{eq:pi_scale}
\end{equation}
where $\delta > 0$ prevents division by zero. The resulting damped inverses are formulated as:
\begin{equation}
\begin{aligned}
\mathbf{A}_{t,\mathrm{inv}}^{(\ell)} &= \left(\hat{\mathbf{A}}^{(\ell)}_t + \sqrt{\lambda^{(t)}}\pi^{(\ell)}\mathbf{I}\right)^{-1}, \\
\mathbf{S}_{t,\mathrm{inv}}^{(\ell)} &= \left(\hat{\mathbf{S}}^{(\ell)}_t + \frac{\sqrt{\lambda^{(t)}}}{\pi^{(\ell)}}\mathbf{I}\right)^{-1}.
\end{aligned}
\end{equation}
Given a combined weight-bias gradient matrix $\mathbf{C}^{(\ell)} = [\mathbf{G}^{(\ell)} \; \mathbf{g}_b^{(\ell)}]$, the curvature-preconditioned gradient is evaluated layer-wise via synchronous matrix multiplication:
\begin{equation}
\tilde{\mathbf{C}}^{(\ell)} = \mathbf{S}_{t,\mathrm{inv}}^{(\ell)} \, \mathbf{C}^{(\ell)} \, \mathbf{A}_{t,\mathrm{inv}}^{(\ell)}.
\label{eq:kfac_precond_op}
\end{equation}

\subsection{Second-Difference Curvature Sensing}
\label{subsec:rdsa_bg}
Rather than estimating the gradient, symmetric directional evaluations can be used to construct a second-difference identity that measures local curvature. For a parameter block $\boldsymbol{\theta}^{(\ell)}$, let $\mathbf{z} \sim \mathcal{N}(\mathbf{0}, \mathbf{I}_{d_\ell})$ be a random Gaussian perturbation. By evaluating the loss at $\boldsymbol{\theta}^{(\ell)} \pm \epsilon_{\mathrm{Z}} \mathbf{z}$, the second difference is defined as:
\begin{equation}
c = \frac{\mathcal{L}(\boldsymbol{\theta}^{(\ell)} + \epsilon_{\mathrm{Z}} \mathbf{z}; \mathcal{B}) - 2\mathcal{L}(\boldsymbol{\theta}^{(\ell)}; \mathcal{B}) + \mathcal{L}(\boldsymbol{\theta}^{(\ell)} - \epsilon_{\mathrm{Z}} \mathbf{z}; \mathcal{B})}{\epsilon_{\mathrm{Z}}^2}.
\label{eq:second_diff}
\end{equation}
By Taylor expansion, this scalar $c$ is an unbiased estimator of the quadratic form $\mathbf{z}^T \mathbf{H}^{(\ell)} \mathbf{z}$, where $\mathbf{H}^{(\ell)}$ is the local Hessian of the loss with respect to block $\ell$.

% \subsection{Pipeline Foundations and Evaluation Metrics}
% \label{subsec:pipeline_metrics}

% \paragraph{Regularized Training Objectives.} To minimize model over-fitting and improve final generalization, the optimization objectives are wrapped in modern regularization functions. Label smoothing transforms sparse categorical targets into a distributed probability space using a smoothing index $\varepsilon_{\mathrm{ls}} \in (0, 1)$, adjusting the target vector $\mathbf{q} \in \mathbb{R}^C$ such that $q_k = 1 - \varepsilon_{\mathrm{ls}} + \frac{\varepsilon_{\mathrm{ls}}}{C}$ for the correct class, and $q_k = \frac{\varepsilon_{\mathrm{ls}}}{C}$ otherwise. Additionally, Mixup training forces linear behavior between samples by constructing convex combinations of pairs of inputs $(\mathbf{x}_i, \mathbf{x}_j)$ via mixing weights drawn from a Beta distribution:
% \begin{equation}
% \tilde{\mathbf{x}} = \lambda \mathbf{x}_i + (1-\lambda)\mathbf{x}_j, \qquad \lambda \sim \mathrm{Beta}(\alpha, \alpha).
% \label{eq:mixup_blend}
% \end{equation}
% The active optimization loss under Mixup interpolation is consequently evaluated as:
% \begin{equation}
% \ell_{\mathrm{mix}}(\tilde{\mathbf{x}}; \boldsymbol{\theta}) = \lambda\,\ell_{\mathrm{ls}}(\tilde{\mathbf{x}}, y_i; \boldsymbol{\theta}) + (1-\lambda)\,\ell_{\mathrm{ls}}(\tilde{\mathbf{x}}, y_j; \boldsymbol{\theta}).
% \label{eq:mixup_loss_final}
% \end{equation}

% \paragraph{Evaluation Metrics.} To comprehensively measure optimization trajectory dynamics and test generalization performance, we evaluate four key metrics:
% 1. \textbf{Top-1 and Top-5 Accuracy ($\mathrm{Acc@1}, \mathrm{Acc@5}$):} Measures standard categorical classification performance.
% 2. \textbf{Generalization Gap ($\Delta_{\mathrm{gen}}$):} Defined as the signed difference between the running training loss and validation loss at epoch $e$: $\Delta_{\mathrm{gen}}^{(e)} = \mathcal{L}_{\mathrm{train}}^{(e)} - \mathcal{L}_{\mathrm{val}}^{(e)}$.
% 3. \textbf{Epochs-to-Threshold ($e^{\star}(\tau)$):} The minimum training epoch required to cross a predefined target performance metric $\tau$: $e^{\star}(\tau) = \min \{ e : \mathrm{Acc}^{(e)}_{\mathrm{train}} \ge \tau \}$. This serves as our primary proxy for raw empirical convergence velocity.

% This section introduces the notation and the optimization components used in the proposed method. The presentation is restricted to the techniques that are explicitly implemented in the training code: supervised classification with label smoothing, Mixup augmentation, Adam with Gradient Centralization, K-FAC-based curvature preconditioning, SPSA-based zeroth-order gradient augmentation, global gradient clipping, warmup with cosine warm restarts, and the two baselines SGD with Nesterov momentum and AdamW.

% \subsection{Notation}
% \label{subsec:notation}

% Let a neural network with parameters
% \[
% \boldsymbol{\theta} \in \mathbb{R}^{d}
% \]
% define a mapping from an input $\mathbf{x}$ to class logits
% \[
% \mathbf{z}(\mathbf{x}; \boldsymbol{\theta}) \in \mathbb{R}^{C},
% \]
% where $C=100$ for CIFAR-100. The predicted class probabilities are obtained by softmax:
% \begin{equation}
% p_k(\mathbf{x}; \boldsymbol{\theta})
% =
% \frac{\exp(z_k(\mathbf{x}; \boldsymbol{\theta}))}
% {\sum_{j=1}^{C} \exp(z_j(\mathbf{x}; \boldsymbol{\theta}))},
% \qquad k=1,\dots,C.
% \label{eq:softmax}
% \end{equation}

% For a mini-batch $\mathcal{B} = \{(\mathbf{x}_i,y_i)\}_{i=1}^{B}$, the training loss is written as
% \begin{equation}
% \mathcal{L}(\boldsymbol{\theta}; \mathcal{B})
% =
% \frac{1}{B}\sum_{i=1}^{B}\ell(\mathbf{x}_i,y_i;\boldsymbol{\theta}).
% \label{eq:batch_loss}
% \end{equation}

% At optimization step $t$, the gradient of the batch loss is
% \begin{equation}
% \mathbf{g}^{(t)} = \nabla_{\boldsymbol{\theta}} \mathcal{L}(\boldsymbol{\theta}^{(t)};\mathcal{B}^{(t)}).
% \label{eq:full_grad}
% \end{equation}

% For layer $\ell$, let $\mathbf{W}^{(\ell)}$ and $\mathbf{b}^{(\ell)}$ denote the weight and bias parameters. We use $\mathbf{G}^{(\ell)}$ for the corresponding weight gradient, and we denote the layer input activations and output backpropagated sensitivities by $\mathbf{A}^{(\ell)}$ and $\mathbf{S}^{(\ell)}$, respectively. These are the sufficient statistics used by K-FAC.

% \subsection{Label Smoothing Cross-Entropy}
% \label{subsec:label_smoothing}

% The code uses cross-entropy loss with label smoothing parameter $\varepsilon_{\mathrm{ls}} = 0.1$. Label smoothing replaces the hard one-hot target with a softened target distribution, which is standard in modern image classification training. [web:109]

% For a sample with ground-truth class $y$, the smoothed target distribution $\mathbf{q}\in\mathbb{R}^{C}$ is
% \begin{equation}
% q_k=
% \begin{cases}
% 1-\varepsilon_{\mathrm{ls}}+\dfrac{\varepsilon_{\mathrm{ls}}}{C}, & k=y, \\[6pt]
% \dfrac{\varepsilon_{\mathrm{ls}}}{C}, & k\neq y.
% \end{cases}
% \label{eq:label_smoothing}
% \end{equation}

% The per-sample loss is then
% \begin{equation}
% \ell_{\mathrm{ls}}(\mathbf{x},y;\boldsymbol{\theta})
% =
% -\sum_{k=1}^{C} q_k \log p_k(\mathbf{x};\boldsymbol{\theta}).
% \label{eq:ls_ce}
% \end{equation}

% In our setting, $C=100$ and $\varepsilon_{\mathrm{ls}}=0.1$.

% \subsection{Mixup Training}
% \label{subsec:mixup}

% The training pipeline applies Mixup with Beta parameter $0.2$, exactly as implemented in the code. Mixup forms convex combinations of inputs and labels and is known to improve generalization on datasets including CIFAR-100.

% Given two samples $(\mathbf{x}_i,y_i)$ and $(\mathbf{x}_j,y_j)$, Mixup constructs
% \begin{align}
% \tilde{\mathbf{x}} &= \lambda \mathbf{x}_i + (1-\lambda)\mathbf{x}_j, \label{eq:mixup_x} \\
% \tilde{\mathbf{y}} &= \lambda \mathbf{y}_i + (1-\lambda)\mathbf{y}_j, \label{eq:mixup_y}
% \end{align}
% where
% \begin{equation}
% \lambda \sim \mathrm{Beta}(\alpha,\alpha), \qquad \alpha=0.2.
% \label{eq:mixup_lambda}
% \end{equation}

% Because the implementation mixes targets through the loss rather than explicitly constructing $\tilde{\mathbf{y}}$, the effective Mixup loss is
% \begin{equation}
% \ell_{\mathrm{mix}}(\tilde{\mathbf{x}};\boldsymbol{\theta})
% =
% \lambda\,\ell_{\mathrm{ls}}(\tilde{\mathbf{x}},y_i;\boldsymbol{\theta})
% +
% (1-\lambda)\,\ell_{\mathrm{ls}}(\tilde{\mathbf{x}},y_j;\boldsymbol{\theta}).
% \label{eq:mixup_loss}
% \end{equation}

% The batch objective is the mean of \eqref{eq:mixup_loss} over the mini-batch.

% \subsection{Adam First-Order Update}
% \label{subsec:adam}

% The first-order backbone of the proposed optimizer is Adam. Adam maintains exponentially decayed first and second moments of the gradient and uses them to form adaptive parameter-wise steps.

% Let $\mathbf{g}^{(t)}$ denote the gradient used at step $t$. Adam computes
% \begin{align}
% \mathbf{m}^{(t)} &= \beta_1 \mathbf{m}^{(t-1)} + (1-\beta_1)\mathbf{g}^{(t)}, \label{eq:adam_m} \\
% \mathbf{v}^{(t)} &= \beta_2 \mathbf{v}^{(t-1)} + (1-\beta_2)\bigl(\mathbf{g}^{(t)} \odot \mathbf{g}^{(t)}\bigr), \label{eq:adam_v}
% \end{align}
% with bias correction
% \begin{align}
% \hat{\mathbf{m}}^{(t)} &= \frac{\mathbf{m}^{(t)}}{1-\beta_1^t}, \label{eq:adam_mhat} \\
% \hat{\mathbf{v}}^{(t)} &= \frac{\mathbf{v}^{(t)}}{1-\beta_2^t}. \label{eq:adam_vhat}
% \end{align}

% The parameter update is
% \begin{equation}
% \boldsymbol{\theta}^{(t+1)}
% =
% \boldsymbol{\theta}^{(t)}
% -
% \eta^{(t)}
% \frac{\hat{\mathbf{m}}^{(t)}}{\sqrt{\hat{\mathbf{v}}^{(t)}}+\epsilon},
% \label{eq:adam_update}
% \end{equation}
% where in the code
% \[
% \beta_1=0.9,\qquad \beta_2=0.999,\qquad \epsilon=10^{-8}.
% \]

% In the hybrid optimizer, Adam is applied after zeroth-order augmentation and K-FAC preconditioning, using the resulting gradient surrogate.

% \subsection{Gradient Centralization}
% \label{subsec:gc}

% Gradient Centralization (GC) is applied before the Adam moment updates for tensors with dimension at least two, which matches the implementation for convolutional and linear weights. GC centralizes gradients by subtracting their mean over all non-output dimensions. This operation has been proposed as a simple optimization technique that can improve training stability and generalization. [web:80]

% For a weight tensor $\mathbf{W}\in\mathbb{R}^{n\times d_1\times\cdots\times d_r}$ with gradient $\mathbf{G}$, the centralized gradient is
% \begin{equation}
% \tilde{\mathbf{G}}
% =
% \mathbf{G}
% -
% \mathrm{mean}\!\left(\mathbf{G},
% \text{dims}=\{2,\dots,r+1\},
% \text{keepdim}=1\right).
% \label{eq:gc}
% \end{equation}

% In elementwise notation,
% \begin{equation}
% \tilde{G}_{i,j_1,\dots,j_r}
% =
% G_{i,j_1,\dots,j_r}
% -
% \frac{1}{d_1\cdots d_r}
% \sum_{k_1=1}^{d_1}\cdots\sum_{k_r=1}^{d_r}
% G_{i,k_1,\dots,k_r}.
% \label{eq:gc_elementwise}
% \end{equation}

% The hybrid optimizer uses $\tilde{\mathbf{G}}$ in place of $\mathbf{G}$ in the Adam update.

% \subsection{Natural Gradient and K-FAC}
% \label{subsec:kfac}

% The motivation for K-FAC comes from natural gradient descent, where the gradient is preconditioned by the inverse Fisher information matrix. [web:74]

% Let $\mathbf{F}(\boldsymbol{\theta})$ denote the Fisher information matrix. The natural-gradient update is
% \begin{equation}
% \Delta \boldsymbol{\theta}_{\mathrm{nat}}
% =
% -\eta \mathbf{F}(\boldsymbol{\theta})^{-1}
% \nabla_{\boldsymbol{\theta}} \mathcal{L}(\boldsymbol{\theta}).
% \label{eq:natgrad}
% \end{equation}

% Directly computing and inverting $\mathbf{F}$ is intractable for modern neural networks. K-FAC approximates the Fisher block of each layer as a Kronecker product of two smaller factors. For a fully connected layer, the Fisher block for the weight matrix is approximated as
% \begin{equation}
% \mathbf{F}^{(\ell)}
% \approx
% \mathbf{S}^{(\ell)} \otimes \mathbf{A}^{(\ell)},
% \label{eq:kfac_factorization}
% \end{equation}
% where $\mathbf{A}^{(\ell)}$ is the covariance of the layer inputs and $\mathbf{S}^{(\ell)}$ is the covariance of the backpropagated pre-activation derivatives. [web:74]

% Using empirical mini-batch estimates, the code forms
% \begin{align}
% \mathbf{A}^{(\ell)}_{\text{emp}}
% &=
% \frac{1}{N_A}
% \mathbf{X}^{(\ell)\top}\mathbf{X}^{(\ell)},
% \label{eq:emp_A}
% \\
% \mathbf{S}^{(\ell)}_{\text{emp}}
% &=
% \frac{1}{N_S}
% \mathbf{Y}^{(\ell)\top}\mathbf{Y}^{(\ell)},
% \label{eq:emp_S}
% \end{align}
% where $\mathbf{X}^{(\ell)}$ denotes the matrix of unfolded inputs to the layer and $\mathbf{Y}^{(\ell)}$ denotes the matrix of unfolded output sensitivities. For convolutional layers, the code uses an unfold/im2col transformation so that the same matrix formula applies to \texttt{Conv2d} as to \texttt{Linear} layers.

% The running K-FAC factors are updated by exponential moving average with decay $\rho=0.95$:
% \begin{align}
% \hat{\mathbf{A}}^{(\ell)}_t
% &=
% \rho \hat{\mathbf{A}}^{(\ell)}_{t-1}
% +
% (1-\rho)\mathbf{A}^{(\ell)}_{\text{emp},t},
% \label{eq:kfac_running_A}
% \\
% \hat{\mathbf{S}}^{(\ell)}_t
% &=
% \rho \hat{\mathbf{S}}^{(\ell)}_{t-1}
% +
% (1-\rho)\mathbf{S}^{(\ell)}_{\text{emp},t}.
% \label{eq:kfac_running_S}
% \end{align}

% These updates are performed every
% \[
% \tau_{\mathrm{K}}=3
% \]
% optimizer steps in the modified code.

% \subsection{Damped K-FAC Inverse and \texorpdfstring{$\pi$}{pi}-Scaling}
% \label{subsec:kfac_damping}

% To stabilize the inverse, K-FAC uses damped Kronecker factors. The implementation follows the common trace-based balancing idea by defining
% \begin{equation}
% \pi^{(\ell)}
% =
% \sqrt{
% \frac{
% \mathrm{tr}(\hat{\mathbf{A}}^{(\ell)}_t)\,\dim(\hat{\mathbf{S}}^{(\ell)}_t)
% }{
% \mathrm{tr}(\hat{\mathbf{S}}^{(\ell)}_t)\,\dim(\hat{\mathbf{A}}^{(\ell)}_t)+10^{-8}
% }
% }.
% \label{eq:pi_scaling}
% \end{equation}

% The damped inverses are then
% \begin{align}
% \mathbf{A}_{t,\mathrm{inv}}^{(\ell)}
% &=
% \left(
% \hat{\mathbf{A}}^{(\ell)}_t
% +
% \sqrt{\lambda^{(t)}}\,\pi^{(\ell)}\mathbf{I}
% \right)^{-1},
% \label{eq:A_inv}
% \\
% \mathbf{S}_{t,\mathrm{inv}}^{(\ell)}
% &=
% \left(
% \hat{\mathbf{S}}^{(\ell)}_t
% +
% \frac{\sqrt{\lambda^{(t)}}}{\pi^{(\ell)}}\mathbf{I}
% \right)^{-1}.
% \label{eq:S_inv}
% \end{align}

% The damping coefficient is not constant. In the code it follows a cosine annealing schedule over total optimizer steps $T$:
% \begin{equation}
% \lambda^{(t)}
% =
% \lambda_{\infty}
% +
% \left(\lambda_{0}-\lambda_{\infty}\right)
% \frac{1+\cos(\pi t/T)}{2},
% \qquad 0 \le t \le T.
% \label{eq:damping_schedule}
% \end{equation}

% Here $\lambda_0$ is the initial damping and $\lambda_{\infty}$ is the final damping. This matches the implementation variables \texttt{damping} and \texttt{damping\_final}.

% \subsection{K-FAC Preconditioned Gradient}
% \label{subsec:kfac_precondition}

% Let $\mathbf{G}^{(\ell)}$ be the gradient of a layer weight matrix and let $\mathbf{g}_b^{(\ell)}$ be the corresponding bias gradient if a bias is present. The implementation concatenates them:
% \begin{equation}
% \mathbf{C}^{(\ell)}
% =
% \left[
% \mathbf{G}^{(\ell)} \;\; \mathbf{g}_b^{(\ell)}
% \right].
% \label{eq:concat_grad}
% \end{equation}

% The K-FAC preconditioned gradient is then
% \begin{equation}
% \tilde{\mathbf{C}}^{(\ell)}
% =
% \mathbf{S}_{t,\mathrm{inv}}^{(\ell)}
% \,
% \mathbf{C}^{(\ell)}
% \,
% \mathbf{A}_{t,\mathrm{inv}}^{(\ell)}.
% \label{eq:kfac_precond}
% \end{equation}

% The preconditioned weight and bias gradients are extracted from $\tilde{\mathbf{C}}^{(\ell)}$ and passed to the Adam update. For layers without K-FAC support, the code falls back to Adam with Gradient Centralization only.

% \subsection{Zeroth-Order Gradient Augmentation with SPSA}
% \label{subsec:spsa}

% The zeroth-order component of the hybrid optimizer uses two-point simultaneous perturbation stochastic approximation (SPSA). SPSA estimates directional gradient information from two function evaluations, independent of parameter dimension. [web:82]

% At step $t$, draw a random perturbation vector
% \begin{equation}
% \mathbf{z}^{(t)} \sim \mathcal{N}(\mathbf{0},\mathbf{I}).
% \label{eq:spsa_noise}
% \end{equation}

% Define perturbed parameters
% \begin{align}
% \boldsymbol{\theta}^{(t)}_{+}
% &=
% \boldsymbol{\theta}^{(t)} + \epsilon_{\mathrm{Z}} \mathbf{z}^{(t)},
% \label{eq:theta_plus}
% \\
% \boldsymbol{\theta}^{(t)}_{-}
% &=
% \boldsymbol{\theta}^{(t)} - \epsilon_{\mathrm{Z}} \mathbf{z}^{(t)},
% \label{eq:theta_minus}
% \end{align}
% where in the code
% \[
% \epsilon_{\mathrm{Z}} = 10^{-3}.
% \]

% Using two closure evaluations, the scalar directional estimate is
% \begin{equation}
% s^{(t)}
% =
% \frac{
% \mathcal{L}(\boldsymbol{\theta}^{(t)}_{+};\mathcal{B}^{(t)})
% -
% \mathcal{L}(\boldsymbol{\theta}^{(t)}_{-};\mathcal{B}^{(t)})
% }{
% 2\epsilon_{\mathrm{Z}}
% }.
% \label{eq:spsa_scalar}
% \end{equation}

% The perturbation-based gradient contribution is
% \begin{equation}
% \mathbf{g}_{\mathrm{ZO}}^{(t)}
% =
% s^{(t)} \mathbf{z}^{(t)}.
% \label{eq:spsa_vector}
% \end{equation}

% The implementation then augments the current gradient as
% \begin{equation}
% \mathbf{g}^{(t)}
% \leftarrow
% \mathbf{g}^{(t)}
% +
% \frac{\zeta}{\sqrt{d}}
% \mathbf{g}_{\mathrm{ZO}}^{(t)},
% \label{eq:spsa_augment}
% \end{equation}
% where $\zeta$ is the tunable \texttt{zo\_scale} parameter and $d$ is the parameter count of the perturbed tensor. The $\sqrt{d}$ normalization appears explicitly in the code and prevents the zeroth-order term from dominating large tensors.

% This zeroth-order augmentation is applied every
% \[
% \tau_{\mathrm{Z}}=5
% \]
% optimizer steps in the modified implementation.

% \subsection{Global Gradient Clipping}
% \label{subsec:grad_clip}

% Before the optimizer step, gradients are clipped by global norm with threshold
% \[
% c_{\max}=1.0.
% \]
% If the total norm is
% \begin{equation}
% \|\mathbf{g}^{(t)}\|_2
% =
% \left(
% \sum_{j=1}^{d} (g_j^{(t)})^2
% \right)^{1/2},
% \label{eq:grad_norm}
% \end{equation}
% the clipped gradient is
% \begin{equation}
% \bar{\mathbf{g}}^{(t)}
% =
% \mathbf{g}^{(t)}
% \cdot
% \min\left(1,\frac{c_{\max}}{\|\mathbf{g}^{(t)}\|_2}\right).
% \label{eq:grad_clip}
% \end{equation}

% This clipping is applied uniformly across all optimizers in the code.

% \subsection{Warmup and Cosine Warm Restarts}
% \label{subsec:lr_schedule}

% All compared optimizers use the same learning-rate schedule in the modified code. This schedule consists of a linear warmup phase followed by cosine annealing with warm restarts, following the SGDR framework. [web:98]

% Let $\eta_{\mathrm{base}}$ be the optimizer’s tuned base learning rate and let $E_{\mathrm{warm}}=5$ be the number of warmup epochs. During warmup,
% \begin{equation}
% \eta_e
% =
% \eta_{\mathrm{base}}
% \frac{e}{E_{\mathrm{warm}}},
% \qquad e=1,\dots,E_{\mathrm{warm}}.
% \label{eq:warmup}
% \end{equation}

% After warmup, the code uses cosine warm restarts with restart period $T_0=50$, multiplier $T_{\mathrm{mult}}=1$, and minimum learning rate $\eta_{\min}=10^{-5}$. For an epoch index within a restart cycle, the learning rate is
% \begin{equation}
% \eta_e
% =
% \eta_{\min}
% +
% \frac{1}{2}
% (\eta_{\mathrm{base}}-\eta_{\min})
% \left(
% 1+\cos\left(\pi \frac{e_{\mathrm{cycle}}}{T_0}\right)
% \right).
% \label{eq:sgdr}
% \end{equation}

% This common schedule ensures that optimizer comparisons are not confounded by different learning-rate policies.

% \subsection{SGD with Nesterov Momentum}
% \label{subsec:sgd_nesterov}

% One baseline in the code is SGD with Nesterov momentum and momentum coefficient $\mu=0.9$. Nesterov momentum evaluates the gradient after a look-ahead step and is a standard strong SGD baseline. The update can be written as [standard formulation]
% \begin{align}
% \mathbf{v}^{(t+1)}
% &=
% \mu \mathbf{v}^{(t)}
% -
% \eta^{(t)}
% \nabla_{\boldsymbol{\theta}}
% \mathcal{L}\bigl(\boldsymbol{\theta}^{(t)}+\mu \mathbf{v}^{(t)}\bigr),
% \label{eq:nesterov_v}
% \\
% \boldsymbol{\theta}^{(t+1)}
% &=
% \boldsymbol{\theta}^{(t)} + \mathbf{v}^{(t+1)}.
% \label{eq:nesterov_theta}
% \end{align}

% The PyTorch implementation used in the code is \texttt{SGD(..., momentum=0.9, nesterov=True)}.

% \subsection{AdamW Baseline}
% \label{subsec:adamw}

% The second baseline is AdamW, which decouples weight decay from the adaptive gradient update rather than folding it into the gradient itself. This is the standard modern variant of Adam for image classification. [web:94]

% AdamW keeps the same moment estimates as Adam in \eqref{eq:adam_m}--\eqref{eq:adam_vhat}, but updates parameters as
% \begin{equation}
% \boldsymbol{\theta}^{(t+1)}
% =
% \boldsymbol{\theta}^{(t)}
% -
% \eta^{(t)}
% \frac{\hat{\mathbf{m}}^{(t)}}{\sqrt{\hat{\mathbf{v}}^{(t)}}+\epsilon}
% -
% \eta^{(t)}\lambda_{\mathrm{wd}}\boldsymbol{\theta}^{(t)}.
% \label{eq:adamw}
% \end{equation}

% This differs from Adam with $L_2$ regularization, where the decay term is absorbed into the gradient before moment estimation.

% \subsection{Trinity Update Summary}
% \label{subsec:hybrid_summary}

% Combining the above components, the hybrid optimizer used in the code can be summarized as follows.

% First, a backpropagated gradient is computed from the Mixup and label-smoothed classification loss:
% \begin{equation}
% \mathbf{g}_{\mathrm{BP}}^{(t)}
% =
% \nabla_{\boldsymbol{\theta}}
% \mathcal{L}_{\mathrm{mix}}(\boldsymbol{\theta}^{(t)};\mathcal{B}^{(t)}).
% \label{eq:bp_grad}
% \end{equation}

% Second, if $t \bmod \tau_{\mathrm{Z}}=0$, a zeroth-order SPSA term is added:
% \begin{equation}
% \mathbf{g}_{\mathrm{aug}}^{(t)}
% =
% \mathbf{g}_{\mathrm{BP}}^{(t)}
% +
% \frac{\zeta}{\sqrt{d}}\mathbf{g}_{\mathrm{ZO}}^{(t)}.
% \label{eq:aug_grad}
% \end{equation}

% Third, for supported layers, the gradient is preconditioned by K-FAC:
% \begin{equation}
% \mathbf{g}_{\mathrm{pre}}^{(t)}
% =
% \mathcal{P}_{\mathrm{KFAC}}\bigl(\mathbf{g}_{\mathrm{aug}}^{(t)};\lambda^{(t)}\bigr),
% \label{eq:pre_grad}
% \end{equation}
% where $\mathcal{P}_{\mathrm{KFAC}}(\cdot)$ denotes the layerwise operation defined in \eqref{eq:kfac_precond}.

% Fourth, for tensors of dimension at least two, Gradient Centralization is applied:
% \begin{equation}
% \mathbf{g}_{\mathrm{gc}}^{(t)}
% =
% \mathcal{C}\bigl(\mathbf{g}_{\mathrm{pre}}^{(t)}\bigr),
% \label{eq:gc_operator}
% \end{equation}
% where $\mathcal{C}(\cdot)$ denotes the centralization operator in \eqref{eq:gc}.

% Finally, the resulting gradient enters Adam:
% \begin{equation}
% \boldsymbol{\theta}^{(t+1)}
% =
% \mathrm{Adam}\bigl(
% \boldsymbol{\theta}^{(t)},
% \mathbf{g}_{\mathrm{gc}}^{(t)},
% \eta^{(t)}
% \bigr).
% \label{eq:hybrid_final}
% \end{equation}

% This sequence reflects the actual order implemented in the optimizer: zeroth-order augmentation, K-FAC preconditioning, Gradient Centralization, and Adam-style adaptive stepping.

% \subsection{Evaluation Quantities}
% \label{subsec:metrics}

% The code reports the following quantities during training and evaluation.

% \paragraph{Top-1 accuracy.}
% \begin{equation}
% \mathrm{Acc@1}
% =
% \frac{1}{N}
% \sum_{i=1}^{N}
% \mathbb{I}
% \left[
% y_i = \arg\max_{k} p_k(\mathbf{x}_i;\boldsymbol{\theta})
% \right].
% \label{eq:top1}
% \end{equation}

% \paragraph{Top-5 accuracy.}
% \begin{equation}
% \mathrm{Acc@5}
% =
% \frac{1}{N}
% \sum_{i=1}^{N}
% \mathbb{I}
% \left[
% y_i \in \mathrm{Top5}\bigl(\mathbf{p}(\mathbf{x}_i;\boldsymbol{\theta})\bigr)
% \right].
% \label{eq:top5}
% \end{equation}

% \paragraph{Generalization gap.}
% The code logs the signed difference between train loss and validation loss:
% \begin{equation}
% \Delta_{\mathrm{gen}}^{(e)}
% =
% \mathcal{L}_{\mathrm{train}}^{(e)}
% -
% \mathcal{L}_{\mathrm{val}}^{(e)}.
% \label{eq:gen_gap}
% \end{equation}

% \paragraph{Epochs-to-threshold.}
% For a target training accuracy threshold $\tau$, the code records
% \begin{equation}
% e^{\star}(\tau)
% =
% \min\left\{
% e : \mathrm{Acc}^{(e)}_{\mathrm{train}} \ge \tau
% \right\}.
% \label{eq:epochs_to_threshold}
% \end{equation}

% These metrics are used to compare convergence speed, final performance, and generalization behavior across Trinity, SGD with Nesterov momentum, and AdamW.


\section{Proposed Approach}
\label{sec:proposed_method}

Trinity is the systematic orchestration of zeroth, second, and first-order optimization primitives within a unified parameter update pipeline. Rather than viewing these three paradigms as competing methodologies, Trinity treats them as complementary information streams that act on distinct geometric properties of the non-convex loss landscape. This section formalizes the integration mechanics, the non-interfering state tracking, and the structural sequence of operations that guarantee numerical stability. The complete, step-by-step execution of this multi-order pipeline is formally detailed in Algorithm~\ref{alg:trinity}

\subsection{Algorithmic Interleaving and Sequence of Operations}
\label{subsec:sequence}
Trinity processes gradients through a strict, sequential pipeline. At each optimization step $t$, the transition from the raw backpropagated gradient to the final parameter step adheres to the following structural sequence:
\begin{align}
\mathbf{g}_{\mathrm{BP}}^{(t)} &\xrightarrow{\text{ZO Augmentation}} \mathbf{g}_{\mathrm{aug}}^{(t)} \\
&\xrightarrow{\text{SO Preconditioning}} \mathbf{g}_{\mathrm{pre}}^{(t)} \\
&\xrightarrow{\text{Gradient Centralization}} \mathbf{g}_{\mathrm{gc}}^{(t)} \\
&\xrightarrow{\text{First-Order Adam Engine}} \boldsymbol{\theta}^{(t+1)}.
\end{align}

Injecting the zeroth-order exploratory signal \textit{prior} to second-order preconditioning ensures that the exploratory random perturbation is dynamically scaled by the localized curvature tensor.

\subsection{Non-Interfering Zeroth-Order Exploration}
\label{subsec:zo_mechanics}

The zeroth-order component executes at fixed intervals ($\tau_{\mathrm{Z}}$) to provide a localized landscape-aware exploration signal. However, evaluating the loss at perturbed weights $\boldsymbol{\theta}^{(t)} \pm \epsilon_{\mathrm{Z}} \mathbf{z}^{(t)}$ requires two forward passes via a closure function. In a network equipped with K-FAC activation-tracking hooks, these auxiliary forward passes would natively overwrite the forward activation cache ($\mathbf{A}_{\mathrm{raw}}^{(\ell)}$) with out-of-distribution metrics evaluated off the true optimization path.

To solve this issue, Trinity implements a synchronized flag, $\mathcal{F}_{\mathrm{collect}}$. During a baseline optimization step, $\mathcal{F}_{\mathrm{collect}} = \text{True}$, allowing the forward and backward hooks to populate the statistics dictionaries. When $t \equiv 0 \pmod{\tau_{\mathrm{Z}}}$, the optimizer sets $\mathcal{F}_{\mathrm{collect}} = \text{False}$ prior to invoking the closure evaluations:
\section{Proposed Approach: Trinity}
\label{sec:method}

\subsection{Design Principle}
\label{subsec:design_principle}
The core principle of Trinity is a strict reassignment of optimization roles: zeroth-order signals sense the geometry, while first- and second-order methods act on it. A purely additive tri-order composition suffers from $\mathcal{O}(d)$ variance in the ZO gradient estimator and biases the descent direction through layer-wise rescaling. Instead, Trinity structurally decouples the sensing mechanism from the update engine.

\subsection{ZO Curvature Sensor}
\label{subsec:zo_sensor}
Using the second-difference identity (Eq. \ref{eq:second_diff}), Trinity calculates two moment estimators for each parameter block over $m$ probes. For a block $\ell$, the debiased effective rank ratio $\hat{\rho}$ and the negative-curvature fraction $\hat{\nu}$ are computed. Three critical correctness requirements ensure the validity of this sensor:
\begin{enumerate}
    \item \textbf{Shared Minibatch:} The same minibatch must be used across all probes to eliminate variance contamination from data sampling.
    \item \textbf{High-Precision Accumulation:} Probe losses must be accumulated in fp32 or fp64, as the $1/\epsilon_{\mathrm{Z}}^2$ division quickly degrades in bf16/fp16.
    \item \textbf{Perturbation Radius:} A macroscopic radius $\epsilon_{\mathrm{Z}} \approx 10^{-2}$ is required.
\end{enumerate}
To protect the Fisher statistics tracking from the sensor's off-path forward passes, a collection flag ($\mathcal{F}_{\mathrm{collect}}$) gates the forward hooks during sensing.

\subsection{Information-Geometric Interpretation}
\label{subsec:info_geom}
The sensor's effective rank ratio $\hat{\rho}$ carries a direct information-geometric interpretation. As $\hat{\rho} \to 1$, the block's Fisher Information Matrix approaches a scaled identity $\mathbf{F} \approx c\mathbf{I}$. In this regime, the natural gradient degenerates to the Euclidean gradient ($\mathbf{F}^{-1}\nabla L \propto \nabla L$), meaning the expensive $\mathcal{O}(n^3)$ matrix inversion is wasted. Conversely, $\hat{\rho} \ll 1$ indicates highly curved, anisotropic geometry where the Kronecker structure of K-FAC provides significant preconditioning value. This interpretation justifies the use of K-FAC over other preconditioning methods, as the ZO sensor measures the specific geometric object (Hessian/Fisher anisotropy) that the actuator approximates.

\subsection{Round-Robin Scheduling}
\label{subsec:round_robin}
Evaluating the curvature sensor across all $B$ blocks per sensing event would require $(2m+1)B$ forward passes, which is prohibitively expensive. Moreover, per-block probes cannot share forward passes, as a joint perturbation would measure cross-block terms rather than the diagonal block's curvature. Trinity circumvents this by adopting a round-robin schedule over $n_p$ blocks per sensing interval $\tau_{\mathrm{Z}}$. This bounds the sensing overhead to approximately 6\% of total training time, regardless of the architecture size.

\subsection{Controller}
\label{subsec:controller}
The decision to activate K-FAC preconditioning for a given block is governed by a stateful controller. The controller maintains a smoothed estimate $\tilde{\rho}^{(\ell)}$ and employs hysteresis bounds ($\rho_{\mathrm{lo}} < \rho_{\mathrm{hi}}$) to prevent chattering. A minimum dwell time $\tau_{\mathrm{dwell}}$ prevents rapid oscillation between modes. Crucially, a global budget parameter $\beta_{\mathrm{max}}$ restricts the maximum fraction of blocks that can be inverted concurrently. The controller admits the top $\lfloor \beta_{\mathrm{max}} B \rfloor$ candidate blocks with the lowest $\tilde{\rho}$ values. 

\subsection{Grafting}
\label{subsec:grafting}
First-order (FO) and second-order (SO) modes inherently operate on different step-size scales. A naive switch between modes would induce a sudden jump in the effective step size by an unknown, layer-dependent factor. To ensure a scale-safe transition, Trinity employs magnitude grafting. A shadow AdamW state is permanently maintained for every block to track FO moments. When a block is updated, the \textit{direction} is taken from the active mode (SO if preconditioned, FO otherwise), but the \textit{magnitude} (Frobenius norm) is strictly grafted from the shadow AdamW state.

\subsection{ZO as Escape Actuator}
\label{subsec:escape}
In addition to guiding the controller, the ZO sensor functions directly as a saddle-escape actuator. If a block's negative curvature fraction exceeds a threshold ($\tilde{\nu}^{(\ell)} \ge \nu_{\mathrm{thr}}$) while its gradient norm is abnormally small (indicating a saddle point or sharp local minimum), a single perturbed-GD step is executed. A lockout counter $\tau_{\mathrm{esc}}$ prevents immediate re-triggering.

\begin{algorithm*}[!t]
\footnotesize
\caption{Trinity Optimization Framework}
\label{alg:trinity}
\begin{algorithmic}[1]
\Require Parameters $\boldsymbol{\theta}^{(0)}$; learning rate schedule $\eta^{(t)}$; loss function $L(\cdot)$; blocks $\{\ell\}$, sizes $d_\ell$
\Require AdamW: $\beta_1, \beta_2, \epsilon_a, \lambda_{\mathrm{wd}}$
\Require Sensor: $\tau_{\mathrm{Z}}, m, \epsilon_{\mathrm{Z}}, n_p$
\Require Control: $\rho_{\mathrm{lo}}, \rho_{\mathrm{hi}}, \alpha_{\mathrm{ema}}, \tau_{\mathrm{dwell}}, \beta_{\mathrm{max}}$
\Require K-FAC: $\tau_{\mathrm{stat}}, \tau_{\mathrm{inv}}, \rho_{\mathrm{kf}}, \lambda_{\mathrm{damp}}, d_{\mathrm{max}}$
\Require Escape: $\nu_{\mathrm{thr}}, \gamma_g, r_{\mathrm{esc}}, \tau_{\mathrm{esc}}$
\State \textbf{Initialize:} For all $\ell$: $\mathbf{M}, \mathbf{V} \leftarrow 0$ (SO branch); $\mathbf{M}_{\mathrm{ad}}, \mathbf{V}_{\mathrm{ad}} \leftarrow 0$ (shadow AdamW); $\hat{\mathbf{A}}, \hat{\mathbf{S}}, \mathbf{A}_{\mathrm{inv}}, \mathbf{S}_{\mathrm{inv}} \leftarrow \mathbf{I}$
\State \textbf{Initialize:} For all $\ell$: $\mathrm{mode} \leftarrow \mathrm{FO}, \tilde{\rho} \leftarrow 1, \tilde{\nu} \leftarrow 0, \mathrm{dwell} \leftarrow 0, \mathrm{lock} \leftarrow 0, \bar{g} \leftarrow 0$; $\mathcal{F}_{\mathrm{collect}} \leftarrow \mathrm{True}; t \leftarrow 0$
\While{not converged}
    \State $t \leftarrow t + 1$; sample minibatch $\mathcal{B}^{(t)}$
    \State $L_{\mathrm{base}} \leftarrow L(\boldsymbol{\theta}^{(t)}; \mathcal{B}^{(t)})$; $\mathbf{g} \leftarrow \nabla_{\boldsymbol{\theta}} L_{\mathrm{base}}$ \Comment{Hooks cache $\mathbf{A}_{\mathrm{raw}}, \mathbf{S}_{\mathrm{raw}}$}
    \If{use\_GC}
        \State $\mathbf{g}^{(\ell)} \leftarrow \mathrm{CENTRALIZE}(\mathbf{g}^{(\ell)})$ for all $\ell$ with $\mathrm{ndim} \ge 2$
    \EndIf

    \Statex \textbf{Phase 1: ZO Sensing (every $\tau_{\mathrm{Z}}$ steps)}
    \If{$t \bmod \tau_{\mathrm{Z}} == 0$}
        \State $\mathcal{F}_{\mathrm{collect}} \leftarrow \mathrm{False}$ \Comment{Freeze Fisher statistics}
        \State $S_t \leftarrow$ next $n_p$ eligible blocks (round-robin)
        \For{$\ell \in S_t$}
            \State $(\hat{\rho}, \hat{\nu}) \leftarrow \mathrm{ZO\_PROBE}(\boldsymbol{\theta}^{(t)}, \mathcal{B}^{(t)}, \ell, m, \epsilon_{\mathrm{Z}})$ \Comment{Second-difference sensing}
            \State $\tilde{\rho}^{(\ell)} \leftarrow \alpha_{\mathrm{ema}} \tilde{\rho}^{(\ell)} + (1-\alpha_{\mathrm{ema}}) \hat{\rho}$; $\tilde{\nu}^{(\ell)} \leftarrow \hat{\nu}$
        \EndFor
        \State $\mathcal{F}_{\mathrm{collect}} \leftarrow \mathrm{True}$
        \State $\{\mathrm{mode}^{(\ell)}\} \leftarrow \mathrm{CONTROLLER}(\{\tilde{\rho}\}, \{\mathrm{dwell}\}, \beta_{\mathrm{max}})$
    \EndIf

    \Statex \textbf{Phase 2: ZO Actuator (Saddle Escape)}
    \For{$\ell \in S_t$ sensed at this step}
        \If{$\tilde{\nu}^{(\ell)} \ge \nu_{\mathrm{thr}}$ and $\|\mathbf{g}^{(\ell)}\|_2 \le \gamma_g \bar{g}^{(\ell)}$ and $\mathrm{lock}^{(\ell)} == 0$}
            \State $\mathbf{u} \sim \mathrm{Uniform}(\mathcal{S}^{d_\ell - 1})$
            \State $\boldsymbol{\theta}^{(\ell)} \leftarrow \boldsymbol{\theta}^{(\ell)} + r_{\mathrm{esc}} \|\boldsymbol{\theta}^{(\ell)}\|_2 \mathbf{u}$ \Comment{Perturbed GD}
            \State $\mathbf{M}^{(\ell)} \leftarrow \frac{1}{2}\mathbf{M}^{(\ell)}$; $\mathbf{M}_{\mathrm{ad}}^{(\ell)} \leftarrow \frac{1}{2}\mathbf{M}_{\mathrm{ad}}^{(\ell)}$ \Comment{Damp stale momentum}
            \State $\mathrm{lock}^{(\ell)} \leftarrow \tau_{\mathrm{esc}}$
        \EndIf
    \EndFor
    \For{all $\ell$} $\mathrm{lock}^{(\ell)} \leftarrow \max(\mathrm{lock}^{(\ell)} - 1, 0)$; $\bar{g}^{(\ell)} \leftarrow \mathrm{EMA}(\|\mathbf{g}^{(\ell)}\|_2)$ \EndFor

    \Statex \textbf{Phase 3: Fisher Factors and Inverses}
    \If{$t \bmod \tau_{\mathrm{stat}} == 0$}
        \For{all eligible $\ell$}
            \State $\hat{\mathbf{A}}^{(\ell)} \leftarrow \rho_{\mathrm{kf}}\hat{\mathbf{A}}^{(\ell)} + (1-\rho_{\mathrm{kf}})\mathbf{A}_{\mathrm{raw}}^T \mathbf{A}_{\mathrm{raw}}$; $\hat{\mathbf{S}}^{(\ell)} \leftarrow \rho_{\mathrm{kf}}\hat{\mathbf{S}}^{(\ell)} + (1-\rho_{\mathrm{kf}})\mathbf{S}_{\mathrm{raw}}^T \mathbf{S}_{\mathrm{raw}}$
        \EndFor
    \EndIf
    \If{$t \bmod \tau_{\mathrm{inv}} == 0$}
        \For{$\ell$ with $\mathrm{mode}^{(\ell)} == \mathrm{SO}$ and $\dim(\ell) \le d_{\mathrm{max}}$}
            \State $\pi \leftarrow \sqrt{(\mathrm{tr}(\hat{\mathbf{A}}) \cdot \dim(\hat{\mathbf{S}})) / (\mathrm{tr}(\hat{\mathbf{S}}) \cdot \dim(\hat{\mathbf{A}}) + \delta)}$
            \State $\mathbf{A}_{\mathrm{inv}}^{(\ell)} \leftarrow (\hat{\mathbf{A}} + \sqrt{\lambda_{\mathrm{damp}}} \pi \mathbf{I})^{-1}$; $\mathbf{S}_{\mathrm{inv}}^{(\ell)} \leftarrow (\hat{\mathbf{S}} + (\sqrt{\lambda_{\mathrm{damp}}}/\pi) \mathbf{I})^{-1}$
        \EndFor
    \EndIf

    \Statex \textbf{Phase 4: Update Engine with Grafting}
    \For{all $\ell$}
        \State $\mathbf{M}_{\mathrm{ad}} \leftarrow \beta_1 \mathbf{M}_{\mathrm{ad}} + (1-\beta_1)\mathbf{g}^{(\ell)}$; $\mathbf{V}_{\mathrm{ad}} \leftarrow \beta_2 \mathbf{V}_{\mathrm{ad}} + (1-\beta_2)(\mathbf{g}^{(\ell)} \odot \mathbf{g}^{(\ell)})$
        \State $\boldsymbol{\Delta}_{\mathrm{ad}} \leftarrow (\mathbf{M}_{\mathrm{ad}} / (1-\beta_1^t)) \oslash (\sqrt{\mathbf{V}_{\mathrm{ad}} / (1-\beta_2^t)} + \epsilon_a)$ \Comment{Shadow AdamW magnitude}
        \If{$\mathrm{mode}^{(\ell)} == \mathrm{SO}$}
            \State $\mathbf{C} \leftarrow [\mathbf{G}^{(\ell)} \;\; \mathbf{g}_b^{(\ell)}]$
            \State $\mathbf{G}_{\mathrm{pre}}, \mathbf{g}_{b,\mathrm{pre}} \leftarrow \mathrm{Split}(\mathbf{S}_{\mathrm{inv}}^{(\ell)} \mathbf{C} \mathbf{A}_{\mathrm{inv}}^{(\ell)})$
            \State $\mathbf{M} \leftarrow \beta_1 \mathbf{M} + (1-\beta_1)\mathbf{G}_{\mathrm{pre}}$; $\mathbf{V} \leftarrow \beta_2 \mathbf{V} + (1-\beta_2)(\mathbf{G}_{\mathrm{pre}} \odot \mathbf{G}_{\mathrm{pre}})$
            \State $\boldsymbol{\Delta}_{\mathrm{act}} \leftarrow (\mathbf{M} / (1-\beta_1^t)) \oslash (\sqrt{\mathbf{V} / (1-\beta_2^t)} + \epsilon_a)$
        \Else
            \State $\boldsymbol{\Delta}_{\mathrm{act}} \leftarrow \boldsymbol{\Delta}_{\mathrm{ad}}$; $\mathbf{M} \leftarrow \mathbf{M}_{\mathrm{ad}}$; $\mathbf{V} \leftarrow \mathbf{V}_{\mathrm{ad}}$ \Comment{Keep SO state warm}
        \EndIf
        \State $\boldsymbol{\Delta} \leftarrow \|\boldsymbol{\Delta}_{\mathrm{ad}}\|_F \cdot \boldsymbol{\Delta}_{\mathrm{act}} / (\|\boldsymbol{\Delta}_{\mathrm{act}}\|_F + \delta)$ \Comment{GRAFTING}
        \State $\boldsymbol{\theta}^{(\ell)} \leftarrow \boldsymbol{\theta}^{(\ell)} - \eta^{(t)}(\boldsymbol{\Delta} + \lambda_{\mathrm{wd}}\boldsymbol{\theta}^{(\ell)})$
    \EndFor
\EndWhile
\end{algorithmic}
\end{algorithm*}

\section{Experimental Setup}
\label{sec:experimental_setup}

To evaluate the optimization capabilities of the proposed Trinity, we evaluated on image classification benchmark across three distinct structural paradigms: (i) a heavy
convolutional neural network (CNN), (ii) a large-scale vision transformer backbone, and (iii) a lightweight residual network.
To ensure the reproducibility of our claims, the complete PyTorch implementation of the Trinity framework, alongside all benchmarking pipelines, hyperparameter search configurations, and ablation scripts, has been open-sourced in an anonymized repository: \url{https://anonymous.4open.science/r/Trinity-opt-DSAA2026/}
\subsection{Dataset and Imbalance Configuration}
For all evaluations, we utilized a challenging Long-Tailed
version of the CIFAR-10 dataset to simulate real-world label imbalances and out-of-distribution tracking performance. We
applied an architectural imbalance factor of $0.01$, creating a severe
$100:1$ ratio between the dominant majority and rare minority classes. This long-tailed pruning significantly restricted
the training capacity to a total of $12{,}406$ active training samples
across all three experimental settings. Data preprocessing and stochastic augmentations remained strictly aligned across all settings, utilizing standard random cropping ($32\times32$
with a padding of 4) and horizontal flips.

\subsection{Architectural Configurations}
\label{subsec:architectures}

To systematically evaluate the scaling properties and geometric adaptability of our proposed framework across highly diverse parameter regimes, we adopt three structural baselines:

\begin{itemize}
    \item \textbf{Experiment 1 (Wide Residual Network):} We evaluate a Wide-ResNet-101-2 architecture containing approximately $126.9\times 10^6$ parameters. Following standard practice for $32 \times 32$ spatial resolutions, the foundational stem of the network was modified to replace the aggressive $7 \times 7$ strided convolution and max-pooling layers with a $3 \times 3$ stride-1 convolution, preserving spatial fidelity. This model configuration introduces a demanding, parameter-dense regime specifically designed to evaluate the numerical scalability and computational stability of the layer-wise K-FAC second-order preconditioning matrices.
    
    \item \textbf{Experiment 2 (Vision Transformer):} We configure a custom transformer comprising approximately $86\times 10^6$ parameters. The network utilizes a patch tokenization layer executed via a $4\times 4$ convolutional projection (yielding a sequence length of $64$ tokens), a learnable \texttt{cls\_token}, and standard positional embeddings. The structural core consists of $12$ Transformer encoder blocks operating with an embedding dimension of $768$, $12$ independent self-attention heads, and an inner feed-forward network (FFN) dimension of $3072$. A pre-layer-normalization (\texttt{norm\_first}) topology is enforced to stabilize internal gradient dynamics.
    
    \item \textbf{Experiment 3 (ResNet-18):} We instantiate a compact ResNet-18 model containing approximately $11\times 10^6$ parameters. This baseline is initialized completely from scratch without pre-trained weights and is mapped to a 10-class output classification head. It serves as a low-complexity control reference to isolate optimizer efficiency across vastly different architectural regimes.
\end{itemize}

All three architectural configurations are evaluated using an identical data pipeline operating with a batch size of $128$. The parameter paths are regularized via a global gradient norm clipping threshold bounded strictly at $1.0$, and the learning rate trajectories are modulated across training epochs using a \texttt{CosineAnnealingLR} schedule policy.

\subsection{Hyperparameter Tuning}
% We performed precise hyperparameter optimization via the Optuna tuning
% framework over distinct initialization sets to extract peak baseline
% performance. Across all runs, base learning rates ($\eta$) and weight decay factors ($\lambda_{wd}$) were completely tuned, alongside specific scaling constants unique to the Hybrid architecture (K-FAC trace-damping coefficients and SPSA zeroth-order step scales).

% The discovered optimal parameters across the respective model spaces are given as follows.

We performed rigorous hyperparameter optimization via the Optuna tuning framework over continuous configuration spaces to isolate peak baseline performance. Across all baseline models, the base learning rates ($\eta$) and weight decay coefficients ($\lambda_{\mathrm{wd}}$) were systematically optimized. For the Trinity architecture, this search domain was expanded to encapsulate its unique multi-order parameters: specifically, the base Tikhonov damping coefficients ($\lambda$) governing second-order K-FAC inversion stability, and the zeroth-order gradient scaling factors ($\zeta$) controlling Random Direction Stochastic Approximation (RDSA) exploration intensity. The discovered optimal hyperparameter sets across the respective architectural spaces are detailed below.
\smallskip
\noindent\textbf{Wide-ResNet-101-2:}
Trinity ($\eta\!\approx\!0.0037$, $\lambda_{wd}\!\approx\!1.4{\times}10^{-4}$,
damping$\,\approx\!0.0117$, $\tau_{K}\!=\!10$, $\zeta\!\approx\!0.279$)
;
Adam ($\eta\!\approx\!3.4{\times}10^{-4}$,
$\lambda_{wd}\!\approx\!1.3{\times}10^{-4}$) ;
RMSProp ($\eta\!\approx\!4.0{\times}10^{-4}$,
$\lambda_{wd}\!\approx\!1.5{\times}10^{-5}$) ;
SGD ($\eta\!\approx\!0.006$,
$\lambda_{wd}\!\approx\!1.18{\times}10^{-4}$) .

\smallskip
\noindent\textbf{Vision Transformer:}
Trinity ($\eta\!\approx\!2.0{\times}10^{-4}$,
$\lambda_{wd}\!\approx\!2.17{\times}10^{-5}$,
damping$\,\approx\!0.00137$, $\tau_{K}\!=\!5$, $\zeta\!\approx\!0.299$)
;
Adam ($\eta\!\approx\!4.1{\times}10^{-4}$,
$\lambda_{wd}\!\approx\!4.6{\times}10^{-4}$) ;
RMSProp ($\eta\!\approx\!1.1{\times}10^{-4}$,
$\lambda_{wd}\!\approx\!3.93{\times}10^{-5}$) ;
SGD ($\eta\!\approx\!0.0075$, $\lambda_{wd}\!\approx\!9.78{\times}10^{-3}$)
.

\smallskip
\noindent\textbf{ResNet-18:}
Trinity ($\eta\!\approx\!0.00245$,
$\lambda_{wd}\!\approx\!1.60{\times}10^{-4}$,
damping$\,\approx\!0.00674$, $\tau_{K}\!=\!5$, $\zeta\!\approx\!0.462$)
;
Adam ($\eta\!\approx\!2.20{\times}10^{-4}$,
$\lambda_{wd}\!\approx\!5.02{\times}10^{-4}$) ;
RMSProp ($\eta\!\approx\!2.99{\times}10^{-4}$,
$\lambda_{wd}\!\approx\!2.02{\times}10^{-5}$) ;
SGD ($\eta\!\approx\!0.00853$,
$\lambda_{wd}\!\approx\!5.03{\times}10^{-3}$) .

%------------------------------------------------------------------
\section{Results}
\label{sec:results}

We evaluated our proposed framework across all three deep learning paradigms. Tables \ref{tab:highest_results} and \ref{tab:results_exp3} report the validation accuracy and total wall-clock time for each optimizer's final full training run. 

\begin{table}[htbp]
\centering
\caption{Validation Accuracy and Wall-Clock Time on Long-Tailed CIFAR-10 (Exp 1 and Exp 2)}
\label{tab:highest_results}
\setlength{\tabcolsep}{4pt}
\begin{tabular}{lcccc}
\toprule
\multirow{}{}{\textbf{Optimizer}}
  & \multicolumn{2}{c}{\textbf{Exp.\,1: Wide-ResNet}}
  & \multicolumn{2}{c}{\textbf{Exp.\,2: Transformer}} \\
\cmidrule(lr){2-3}\cmidrule(lr){4-5}
  & \textbf{Acc.} & \textbf{Time\,(s)}
  & \textbf{Acc.} & \textbf{Time\,(s)} \\
\midrule
SGD (Nesterov)      & 16.62\% & 414 & 56.52\% & 233 \\
RMSProp             & 34.20\% & 382 & 57.46\% & 234 \\
Adam                & 60.38\% & 385 & 59.25\% & 230 \\
\textbf{Trinity}    & \textbf{63.39\%} & 487 & \textbf{62.96\%} & 312 \\
\bottomrule
\end{tabular}
\end{table}

\begin{table}[htbp]
\centering
\caption{Validation Accuracy and Wall-Clock Time on Long-Tailed CIFAR-10, Exp 3.}
\label{tab:results_exp3}
\setlength{\tabcolsep}{5pt}
\begin{tabular}{lcc}
\toprule
\textbf{Optimizer} & \textbf{Val.\,Acc.} & \textbf{Total Time\,(s)} \\
\midrule
RMSProp             & 27.15\% & 965  \\
SGD (Nesterov)      & 38.96\% & 946  \\
\textbf{Trinity}    & 43.69\% & 1583 \\
Adam                & \textbf{43.88\%} & 969  \\
\bottomrule
\end{tabular}
\end{table}


The empirical findings demonstrate that Trinity consistently establishes a superior performance ceiling in highly complex, parameter-dense architectures, while revealing nuanced structural trade-offs in compact models.

In \textbf{Experiment 1} (Figure \ref{fig:exp1}), training the massive Wide-ResNet architecture from scratch on severely imbalanced data causes vanilla SGD to struggle significantly, plateauing at just 16.62\%. This degradation directly aligns with recent theoretical findings on gradient starvation under heavy-tailed class imbalance \cite{kunstner2024heavy}. Trinity achieves a peak accuracy of 63.39\%, substantially outperforming both Adam (60.38\%) and RMSProp (34.20\%). 

In \textbf{Experiment 2} (Figure \ref{fig:exp2}), the non-local attention mechanics of the Vision Transformer shift the performance boundaries. Even within this heavily normalized sequence space, Trinity achieves a validation accuracy of 62.96\% over 300 epochs, establishing a clear ceiling over Adam (59.25\%), RMSProp (57.46\%), and SGD (56.52\%). Naturally, Trinity incurs a deterministic computational time overhead due to the periodic calculation of dual K-FAC factor inversions and the auxiliary forward passes required by the RDSA estimator (e.g., 487s vs. 385s for Adam in Exp 1, and 312s vs. 230s in Exp 2). However, this explicit structural tracking translates directly to a +3.71 percentage point accuracy gain in the ViT architecture, demonstrating a highly favorable trade-off between step-time complexity and overall geometric convergence capability.

In \textbf{Experiment 3} (Figure \ref{fig:exp3}), the compact ResNet-18 backbone yields a different optimization dynamic. Here, Adam achieves a slightly higher accuracy at 43.88\%, followed by Trinity at 43.69\%. This indicates that for lower-capacity models trained on highly constrained datasets, the aggressive structural preconditioning of Trinity acts as a heavy regularizer. While this prevents overfitting, it limits the model from perfectly memorizing the sparse training distribution compared to unconstrained adaptive tracking.

\definecolor{colorhybrid}{RGB}{31, 119, 180}   % Modern Steel Blue
\definecolor{coloradam}{RGB}{214, 39, 40}     % Academic Muted Red
\definecolor{colorrmsprop}{RGB}{44, 160, 44}  % Balanced Forest Green
\definecolor{colorsgd}{RGB}{255, 127, 14}     % Safety Orange

\pgfplotsset{compat=1.18}

% ──────────────────────────────────────────────────────────────────────────────
% FIGURE 1
% ──────────────────────────────────────────────────────────────────────────────
\begin{figure}[t]
  \centering
  \begin{tikzpicture}
    \begin{axis}[
      width=\columnwidth,
      height=6cm, % Optimized vertical ratio for 2-column layouts
      xlabel={Epoch},
      ylabel={Validation Accuracy (\%)},
      xlabel style={font=\small, yshift=2pt},
      ylabel style={font=\small, xshift=-2pt},
      tick label style={font=\footnotesize},
      xmin=1,   xmax=200,
      ymin=5,   ymax=72,
      xtick={1,25,50,75,100,125,150,175,200},
      ytick={10,20,30,40,50,60,70},
      minor tick num=4,
      grid=both,
      minor grid style={line width=0.2pt, draw=black!5},
      major grid style={line width=0.4pt, draw=black!12, dashed},
      legend pos=south east,
      legend style={
        font=\footnotesize,
        draw=black!15,
        fill=white,
        fill opacity=0.95,
        text opacity=1,
        row sep=2pt,
        inner sep=4pt,
        rounded corners=1pt
      },
      clip=true,
    ]

    % --- Hybrid (Proposed Method) ---
    \addplot[
      line width=1.3pt,
      color=colorhybrid,
      mark=none,
    ] table[x=epoch, y=val_acc, col sep=comma] {exp1_hybrid.csv};
    \addlegendentry{Trinity}

    % --- Adam ---
    \addplot[
      line width=0.9pt,
      color=coloradam,
      mark=none,
    ] table[x=epoch, y=val_acc, col sep=comma] {exp1_adam.csv};
    \addlegendentry{Adam}

    % --- RMSProp ---
    \addplot[
      line width=0.9pt,
      color=colorrmsprop,
      mark=none,
    ] table[x=epoch, y=val_acc, col sep=comma] {exp1_rmsprop.csv};
    \addlegendentry{RMSProp}

    % --- SGD ---
    \addplot[
      line width=0.9pt,
      color=colorsgd,
      mark=none,
    ] table[x=epoch, y=val_acc, col sep=comma] {exp1_sgd.csv};
    \addlegendentry{SGD}

    % --- Clean Callout for Peak Value ---
    \node[
      font=\fontseries{b}\selectfont\footnotesize,
      color=colorhybrid,
      fill=white,
      fill opacity=0.8,
      text opacity=1,
      inner sep=1pt,
      anchor=south east,
    ] at (axis cs:195,64.2) {63.39\%};

    \end{axis}
  \end{tikzpicture}
  \caption{Experiment 1 (WResNet)}
  \label{fig:exp1}
\end{figure}

% ─────────────────────────────────────────────────────────
% FIGURE 2: 300 Epoch Evaluation
% ─────────────────────────────────────────────────────────
\begin{figure}[t]
  \centering
  \begin{tikzpicture}
    \begin{axis}[
      width=\columnwidth,
      height=6cm,
      xlabel={Epoch},
      ylabel={Validation Accuracy (\%)},
      xlabel style={font=\small, yshift=2pt},
      ylabel style={font=\small, xshift=-2pt},
      tick label style={font=\footnotesize},
      xmin=1,   xmax=300,
      ymin=5,   ymax=72,
      xtick={1,50,100,150,200,250,300},
      ytick={10,20,30,40,50,60,70},
      minor tick num=4,
      grid=both,
      minor grid style={line width=0.2pt, draw=black!5},
      major grid style={line width=0.4pt, draw=black!12, dashed},
      legend pos=south east,
      legend style={
        font=\footnotesize,
        draw=black!15,
        fill=white,
        fill opacity=0.95,
        text opacity=1,
        row sep=2pt,
        inner sep=4pt,
        rounded corners=1pt
      },
      clip=true,
    ]

    % --- Hybrid (Proposed Method) ---
    \addplot[
      line width=1.3pt,
      color=colorhybrid,
      mark=none,
    ] table[x=epoch, y=val_acc, col sep=comma] {exp2_hybrid.csv};
    \addlegendentry{Trinity}

    % --- Adam ---
    \addplot[
      line width=0.9pt,
      color=coloradam,
      mark=none,
    ] table[x=epoch, y=val_acc, col sep=comma] {exp2_adam.csv};
    \addlegendentry{Adam}

    % --- RMSProp ---
    \addplot[
      line width=0.9pt,
      color=colorrmsprop,
      mark=none,
    ] table[x=epoch, y=val_acc, col sep=comma] {exp2_rmsprop.csv};
    \addlegendentry{RMSProp}

    % --- SGD ---
    \addplot[
      line width=0.9pt,
      color=colorsgd,
      mark=none,
    ] table[x=epoch, y=val_acc, col sep=comma] {exp2_sgd.csv};
    \addlegendentry{SGD}

    % --- Clean Callout for Peak Value ---
    \node[
      font=\fontseries{b}\selectfont\footnotesize,
      color=colorhybrid,
      fill=white,
      fill opacity=0.8,
      text opacity=1,
      inner sep=1pt,
      anchor=south east,
    ] at (axis cs:295,63.8) {62.96\%};

    \end{axis}
  \end{tikzpicture}
  \caption{Experiment 2 (Transformer)}
  \label{fig:exp2}
\end{figure}

% ─────────────────────────────────────────────────────────
% FIGURE 3
% ─────────────────────────────────────────────────────────
\begin{figure}[t]
  \centering
  \begin{tikzpicture}
    \begin{axis}[
      width=\columnwidth,
      height=6cm,
      xlabel={Epoch},
      ylabel={Validation Accuracy (\%)},
      xlabel style={font=\small, yshift=2pt},
      ylabel style={font=\small, xshift=-2pt},
      tick label style={font=\footnotesize},
      xmin=1,   xmax=85,
      ymin=5,   ymax=50,
      xtick={1,10,20,30,40,50,60,70,85},
      ytick={10,20,30,40,50},
      minor tick num=4,
      grid=both,
      minor grid style={line width=0.2pt, draw=black!5},
      major grid style={line width=0.4pt, draw=black!12, dashed},
      legend pos=south east,
      legend style={
  at={(0.5,-0.18)},
  anchor=north,
  font=\footnotesize,
  draw=black!15,
  fill=white,
  fill opacity=0.95,
  text opacity=1,
  row sep=2pt,
  inner sep=4pt,
  rounded corners=1pt,
  legend columns=4,
},
      clip=true,
    ]

    % --- Trinity (Proposed Method) ---
    \addplot[
      line width=1.3pt,
      color=colorhybrid,
      mark=none,
    ] table[x=epoch, y=val_acc, col sep=comma]
      {exp3_hybrid.csv};
    \addlegendentry{Trinity}

    % --- Adam ---
    \addplot[
      line width=0.9pt,
      color=coloradam,
      mark=none,
    ] table[x=epoch, y=val_acc, col sep=comma]
      {exp3_adam.csv};
    \addlegendentry{Adam}

    % --- RMSProp ---
    \addplot[
      line width=0.9pt,
      color=colorrmsprop,
      mark=none,
    ] table[x=epoch, y=val_acc, col sep=comma]
      {exp3_rmsprop.csv};
    \addlegendentry{RMSProp}

    % --- SGD ---
    \addplot[
      line width=0.9pt,
      color=colorsgd,
      mark=none,
    ] table[x=epoch, y=val_acc, col sep=comma]
      {exp3_sgd.csv};
    \addlegendentry{SGD}

    % --- Peak Trinity callout ---
    \node[
      font=\fontseries{b}\selectfont\footnotesize,
      color=colorhybrid,
      fill=white,
      fill opacity=0.85,
      text opacity=1,
      inner sep=1pt,
      anchor=east,
    ] at (axis cs:80,45.0) {43.69\%};

    % --- Peak Adam callout ---
    \node[
      font=\fontseries{b}\selectfont\footnotesize,
      color=coloradam,
      fill=white,
      fill opacity=0.85,
      text opacity=1,
      inner sep=1pt,
      anchor=east,
    ] at (axis cs:80,41.2) {43.88\%};

    \end{axis}
  \end{tikzpicture}
  \caption{Experiment 3 (ResNet-18)}
  \label{fig:exp3}
\end{figure}
%-----------------------------------------------------
\section{Analysis and Discussion}
\label{sec:analysis}


To study the characteristics of our solver, we ran an ablation study comparing the full iteration mechanism against models stripped of second-order scaling (Trinity\,--\,no\,K-FAC) and models operating without stochastic directional exploration (Trinity\,--\,no\,ZO).
Table~\ref{tab:ablation_comparison} lists the validation values logged over these tracking trajectories across all three architectural
settings.


% \begin{table}[t]
% \centering
% \caption{Ablation Study: Component Performance Ceiling Trajectories
%          (20 training epochs each).}
% \label{tab:ablation_comparison}
% \setlength{\tabcolsep}{4pt}
% \begin{tabular}{lccc}
% \toprule
% \textbf{Variant}
%   & \textbf{Exp.\,1} & \textbf{Exp.\,2} & \textbf{Exp.\,3} \\
%   & \textbf{(WResNet)} & \textbf{(Transformer)} & \textbf{(ResNet-18)} \\
% \midrule
% Trinity\,--\,no\,K-FAC & 45.10\% & 30.54\% & 53.92\% \\
% Trinity\,--\,no\,ZO    & 55.77\% & 28.89\% & 57.57\% \\
% \textbf{Trinity (full)}
%   & \textbf{56.82\%} & \textbf{28.84\%} & \textbf{57.79\%} \\
% \bottomrule
% \end{tabular}
% \end{table}

% \begin{table}[htbp]
% \centering
% \caption{Ablation Study: Component Performance Ceiling Trajectories \textbf{(100 training epochs each)}.}
% \label{tab:ablation_comparison}
% \setlength{\tabcolsep}{4pt}
% \begin{tabular}{lccc}
% \toprule
% \textbf{Variant}
%   & \textbf{Exp.\,1} & \textbf{Exp.\,2} & \textbf{Exp.\,3} \\
%   & \textbf{(WResNet)} & \textbf{(Transformer)} & \textbf{(ResNet-18)} \\
% \midrule
% Adam (FO only)         & 59.16\% & 41.74\% & 55.77\% \\
% RDSA (ZO only)         & 59.64\% & 42.41\% & 56.13\% \\
% K-FAC (SO only)        & 63.31\% & 42.83\% & 56.03\% \\
% Trinity\,--\,no\,K-FAC & 59.47\% & 41.75\% & 55.00\% \\ 
% Trinity\,--\,no\,ZO    & 62.79\% & 42.90\% & 57.17\% \\
% \textbf{Trinity (full)}& \textbf{62.52\%} & \textbf{43.87\%} & \textbf{57.23\%} \\
% \bottomrule
% \end{tabular}
% \end{table}




\begin{table}[htbp]
\centering
\caption{Ablation Study (100 epochs). Format: \textbf{Acc.\% / Loss}}
\label{tab:ablation_comparison}
\setlength{\tabcolsep}{3pt} % tighter spacing
\resizebox{\columnwidth}{!}{%
\begin{tabular}{lccc}
\toprule
\textbf{Variant} & \textbf{Exp 1} & \textbf{Exp 2} & \textbf{Exp 3} \\
& \textbf{(WResNet)} & \textbf{(Transformer)} & \textbf{(ResNet18)} \\
\midrule
Adam (FO)     & 59.16 / 0.010 & 41.74 / 0.635 & 55.77 / 0.013 \\
RDSA (ZO)     & 59.64 / 0.009 & 42.41 / 0.629 & 56.13 / 0.023 \\
K-FAC (SO)    & 63.31 / 0.004 & 42.83 / 0.634 & 56.03 / 0.008 \\
Tri -- no K-FAC & 59.47 / 0.008 & 41.75 / 0.648 & 55.00 / 0.017 \\
Tri -- no ZO    & 62.79 / 0.004 & 42.90 / 0.632 & 57.17 / 0.006 \\
\textbf{Trinity (full)} & \textbf{62.52 / 0.003} & \textbf{43.87 / 0.610} & \textbf{57.23 / 0.007} \\
\bottomrule
\end{tabular}%
}
\end{table}




% \caption{Ablation Study: Component Performance Ceiling Trajectories \textbf{(100 training epochs each)}.}
% \label{tab:ablation_comparison}
% \setlength{\tabcolsep}{4pt}
% \begin{tabular}{lccc}
% \toprule
% \textbf{Variant}
%   & \textbf{Exp.\,1} & \textbf{Exp.\,2} & \textbf{Exp.\,3} \\
%   & \textbf{(WResNet)} & \textbf{(Transformer)} & \textbf{(ResNet-18)} \\
% \midrule
% Adam (FO only)         & - & - & 55.77\% \\
% RDSA (ZO only)         & - & - & 56.13\% \\
% K-FAC (SO only)        & - & - & 56.03\% \\
% Trinity\,--\,no\,K-FAC & 45.67\% & 53.92\% & 55.00\% \\ 
% Trinity\,--\,no\,ZO    & 55.60\% & 57.57\% & 57.17\% \\
% \textbf{Trinity (full)}& \textbf{56.72\%} & 57.79\% & \textbf{57.23\%} \\
% \bottomrule
% \end{tabular}
% \end{table}


% \begin{table}[htbp]
% \centering
% \caption{Ablation Study for ResNet-18 model \textbf{(100 training epochs each)}.}
% \label{tab:ablation_comparison}
% \setlength{\tabcolsep}{4pt}
% \begin{tabular}{lccc}
% \toprule
% \textbf{Variant}
%   &  \textbf{Exp.\,3} \\
%   & \textbf{(ResNet-18)} \\
% \midrule
% Adam (FO only)         & 55.77\% \\
% RDSA (ZO only)         & 56.13\% \\
% K-FAC (SO only)        & 56.03\% \\
% Trinity\,--\,no\,K-FAC & 55.00\% \\ 
% Trinity\,--\,no\,ZO    & 57.17\% \\
% \textbf{Trinity (full)}& \textbf{57.23\%} \\
% \bottomrule
% \end{tabular}
% \end{table}


% The comparative breakdown across architectures yields distinct structural insights. For the Wide-ResNet convolutional framework, second-order preconditioning forms the primary driver of performance; omitting K-FAC causes validation metrics to fall steeply to $45.67\%$. The inclusion of multi-point RDSA perturbations provides a steady boost, moving the accuracy boundary up from $55.60\%$ to $56.72\%$.

% When optimizing the token-based sequence spaces of the Vision Transformer, we observed that excluding K-FAC updates leads to a noticeable degradation in early convergence, dropping the trajectory to $53.92\%$, which underscores the necessity of preconditioning even within highly normalized attention mechanisms. Similarly, disabling RDSA exploration decreases the transformer ceiling to $57.57\%$, emphasizing that periodic directional perturbations are vital for navigating the sharp local minima of ViT landscapes toward wider, more generalizable plateaus. The combined Trinity framework achieves the highest peak at $57.79\%$.

% The ablation results for Experiment 3 (ResNet-18) on CIFAR-100 dataset reveal that K-FAC preconditioning remains necessary for establishing early geometric stability (dropping to $26.62\%$ without it). This confirms that second-order curvature information benefits even compact residual architectures under severe class imbalance. Interestingly, disabling the ZO term yields a marginally higher early-stage checkpoint ($28.38\%$ compared to the full $28.28\%$). This indicates that in highly constrained, lower-capacity regimes, the variance injected by the zeroth-order term acts as a heavy fine-grained regularizer that slightly moderates the initial memorization velocity. However, the full Trinity framework demonstrates robust structural integration, effectively scaling its synergistic geometric properties across all three architectural paradigms.


The comparative breakdown across architectures yields distinct structural insights into the optimizer's components. In Experiment 1, we trained a Wide-ResNet-101-2 model from scratch on the long-tailed CIFAR-10 dataset (imbalance factor 0.01) for 100 epochs using a cosine annealing learning rate schedule (base LR $1\text{e-}3$). For this deep convolutional framework, second-order preconditioning forms the primary driver of performance; omitting K-FAC (Trinity\,--\,no\,K-FAC) causes validation accuracy to fall steeply to $59.47\%$. Interestingly, while pure K-FAC (SO only) achieves a high validation peak ($63.31\%$), the full Trinity framework achieves the absolute lowest final training loss across all variants ($0.003$). This indicates that the inclusion of multi-point RDSA perturbations (applied with a scaling factor of $0.1$) acts as a heavy fine-grained regularizer; it slightly moderates the peak validation accuracy to $62.52\%$ while simultaneously allowing the model to explore deeper, more optimal regions of the training landscape.

In Experiment 2, we optimized a custom 12-layer Vision Transformer (embedding dimension 768, 12 heads, patch size $4\times4$) on the same long-tailed CIFAR-10 dataset (imbalance factor 0.01) for 100 epochs. Within these token-based sequence spaces, the synergistic properties of the full Trinity framework become highly apparent, achieving the highest peak accuracy at $43.87\%$ alongside the lowest training loss ($0.610$). We observed that excluding K-FAC updates leads to a degradation in geometric stability, dropping the trajectory down to $41.75\%$, which underscores the necessity of preconditioning even within highly normalized attention mechanisms. Similarly, disabling RDSA exploration (Trinity\,--\,no\,ZO) decreases the transformer ceiling to $42.90\%$, emphasizing that periodic directional perturbations are vital for navigating the sharp local minima of ViT landscapes toward wider, more generalizable plateaus.

Finally, the ablation results for Experiment 3 evaluated a standard ResNet-18 architecture, trained from scratch for 100 epochs on the long-tailed CIFAR-10 dataset (imbalance factor 0.01). The results reveal that K-FAC preconditioning remains strictly necessary for establishing early geometric stability (dropping to $55.00\%$ without it). This confirms that second-order curvature information (computed every 10 steps with a Tikhonov damping factor of $0.01$) benefits even compact residual architectures under severe class imbalance. Furthermore, disabling the ZO term yields a marginally lower validation ceiling ($57.17\%$ compared to the full $57.23\%$), demonstrating that the variance injected by the zeroth-order term successfully pushes the network out of sub-optimal saddle points. Ultimately, the full Trinity framework demonstrates robust structural integration, effectively scaling its synergistic geometric and exploratory properties to achieve optimal or highly competitive convergence across all three architectural paradigms.


\section{Conclusion}
\label{sec:conclusion}

In this work, we presented \textbf{Trinity}, a unified optimization framework that systematically orchestrates zeroth, first, and second-order optimization update rules into a synchronized, non-interfering execution pipeline. By mathematically interleaving localized finite-difference exploration via multi-point RDSA with structural K-FAC preconditioning and Gradient Centralization, our framework takes a step toward exploring the possibilities of adaptivity in optimization methods in terms of stochastic gradient dynamics, information geometric signals, and adaptive moment tracking. Extensive empirical evaluation across a number of parametric scales, spanning deep CNNs and vision transformers, under extreme long-tailed label imbalances, demonstrates that Trinity outperforms the first-order state-of-the-art baselines in many cases. Our systematic ablation analysis validates that the synthesis of randomized gradient augmentation and localized curvature adaptation yields non-redundant advantages, underscoring the profound potential of multi-order hybrid paradigms for navigating deeply non-convex, highly ill-conditioned optimization landscapes in modern deep learning architectures. 

\bibliographystyle{plain}
\begin{thebibliography}{12}

\bibitem{kingma2014adam}
D.~P. Kingma and J.~Ba, ``Adam: A method for stochastic optimization,'' in \emph{International Conference on Learning Representations (ICLR)}, 2015.

\bibitem{loshchilov2019decoupled}
I.~Loshchilov and F.~Hutter, ``Decoupled weight decay regularization,'' in \emph{International Conference on Learning Representations (ICLR)}, 2019.

\bibitem{yong2020gradient}
H.~Yong, J.~Huang, X.~Hua, and L.~Zhang, ``Gradient centralization: A new optimization technique for deep neural networks,'' in \emph{European Conference on Computer Vision (ECCV)}, 2020, pages 635--652.

\bibitem{amari1998natural}
S.-I. Amari, ``Natural gradient works efficiently in learning,'' \emph{Neural Computation}, vol. 10, no. 2, pages 251--276, 1998.

\bibitem{martens2015optimizing}
J.~Martens and R.~Grosse, ``Optimizing neural networks with Kronecker-factored approximate curvature,'' in \emph{International Conference on Machine Learning (ICML)}, 2015, pages 2408--2417.

\bibitem{osawa2019large}
K.~Osawa, Y.~Tsuji, Y.~Ueno, A.~Naruse, R.~Yokota, and S.~Matsuoka, ``Large-scale distributed second-order optimization using Kronecker-factored approximate curvature for deep convolutional neural networks,'' in \emph{Proceedings of the IEEE/CVF Conference on Computer Vision and Pattern Recognition (CVPR)}, 2019.

\bibitem{spall1992multivariate}
J.~C. Spall, ``Multivariate stochastic approximation using a simultaneous perturbation gradient approximation,'' \emph{IEEE Transactions on Automatic Control}, vol. 37, no. 3, pages 332--341, 1992.

\bibitem{nesterov2017random}
Y.~Nesterov and V.~Spokoiny, ``Random gradient-free minimization of convex functions,'' \emph{Foundations of Computational Mathematics}, vol. 17, no. 2, pp. 527--566, 2017.
% \bibitem{foret2020sharpness}
% P.~Foret, A.~Kleiner, H.~Mobahi, and B.~Neyshabur, ``Sharpness-aware minimization for efficiently improving generalization,'' in \emph{International Conference on Learning Representations (ICLR)}, 2021.

% \bibitem{kwon2021asam}
% J.~Kwon, J.~Kim, H.~Park, and I.~K. Choi, ``ASAM: Adaptive sharpness-aware minimization for scale-invariant learning of deep neural networks,'' in \emph{International Conference on Machine Learning (ICML)}, 2021, pages 5905--5914.

\bibitem{zhang2019lookahead}
M.~Zhang, J.~Lucas, J.~Ba, and G.~E. Hinton, ``Lookahead optimizer: k steps forward, 1 step back,'' in \emph{Advances in Neural Information Processing Systems (NeurIPS)}, 2019.

\bibitem{sst2}
R.~Socher, A.~Perelygin, J.~Wu, J.~Chuang, C.~D. Manning, A.~Ng, and C.~Potts, ``Recursive deep models for semantic compositionality over a sentiment treebank,'' in \emph{Proceedings of the Conference on Empirical Methods in Natural Language Processing (EMNLP)}, 2013.

\bibitem{kunstner2024heavy}
F.~Kunstner, R.~Yadav, A.~Milligan, M.~Schmidt, and A.~Bietti, ``Heavy-Tailed Class Imbalance and Why Adam Outperforms Gradient Descent on Language Models,'' \emph{arXiv preprint arXiv:2402.19449}, 2024.


\end{thebibliography}









% \section{Introduction}
% This document is a model and instructions for \LaTeX.
% Please observe the conference page limits. 

% \section{Ease of Use}

% \subsection{Maintaining the Integrity of the Specifications}

% The IEEEtran class file is used to format your paper and style the text. All margins, 
% column widths, line spaces, and text fonts are prescribed; please do not 
% alter them. You may note peculiarities. For example, the head margin
% measures proportionately more than is customary. This measurement 
% and others are deliberate, using specifications that anticipate your paper 
% as one part of the entire proceedings, and not as an independent document. 
% Please do not revise any of the current designations.

% \section{Prepare Your Paper Before Styling}
% Before you begin to format your paper, first write and save the content as a 
% separate text file. Complete all content and organizational editing before 
% formatting. Please note sections \ref{AA}--\ref{SCM} below for more information on 
% proofreading, spelling and grammar.

% Keep your text and graphic files separate until after the text has been 
% formatted and styled. Do not number text heads---{\LaTeX} will do that 
% for you.

% \subsection{Abbreviations and Acronyms}\label{AA}
% Define abbreviations and acronyms the first time they are used in the text, 
% even after they have been defined in the abstract. Abbreviations such as 
% IEEE, SI, MKS, CGS, ac, dc, and rms do not have to be defined. Do not use 
% abbreviations in the title or heads unless they are unavoidable.

% \subsection{Units}
% \begin{itemize}
% \item Use either SI (MKS) or CGS as primary units. (SI units are encouraged.) English units may be used as secondary units (in parentheses). An exception would be the use of English units as identifiers in trade, such as ``3.5-inch disk drive''.
% \item Avoid combining SI and CGS units, such as current in amperes and magnetic field in oersteds. This often leads to confusion because equations do not balance dimensionally. If you must use mixed units, clearly state the units for each quantity that you use in an equation.
% \item Do not mix complete spellings and abbreviations of units: ``Wb/m\textsuperscript{2}'' or ``webers per square meter'', not ``webers/m\textsuperscript{2}''. Spell out units when they appear in text: ``. . . a few henries'', not ``. . . a few H''.
% \item Use a zero before decimal points: ``0.25'', not ``.25''. Use ``cm\textsuperscript{3}'', not ``cc''.)
% \end{itemize}

% \subsection{Equations}
% Number equations consecutively. To make your 
% equations more compact, you may use the solidus (~/~), the exp function, or 
% appropriate exponents. Italicize Roman symbols for quantities and variables, 
% but not Greek symbols. Use a long dash rather than a hyphen for a minus 
% sign. Punctuate equations with commas or periods when they are part of a 
% sentence, as in:
% \begin{equation}
% a+b=\gamma\label{eq}
% \end{equation}

% Be sure that the 
% symbols in your equation have been defined before or immediately following 
% the equation. Use ``\eqref{eq}'', not ``Eq.~\eqref{eq}'' or ``equation \eqref{eq}'', except at 
% the beginning of a sentence: ``Equation \eqref{eq} is . . .''

% \subsection{\LaTeX-Specific Advice}

% Please use ``soft'' (e.g., \verb|\eqref{Eq}|) cross references instead
% of ``hard'' references (e.g., \verb|(1)|). That will make it possible
% to combine sections, add equations, or change the order of figures or
% citations without having to go through the file line by line.

% Please don't use the \verb|{eqnarray}| equation environment. Use
% \verb|{align}| or \verb|{IEEEeqnarray}| instead. The \verb|{eqnarray}|
% environment leaves unsightly spaces around relation symbols.

% Please note that the \verb|{subequations}| environment in {\LaTeX}
% will increment the main equation counter even when there are no
% equation numbers displayed. If you forget that, you might write an
% article in which the equation numbers skip from (17) to (20), causing
% the copy editors to wonder if you've discovered a new method of
% counting.

% {\BibTeX} does not work by magic. It doesn't get the bibliographic
% data from thin air but from .bib files. If you use {\BibTeX} to produce a
% bibliography you must send the .bib files. 

% {\LaTeX} can't read your mind. If you assign the same label to a
% subsubsection and a table, you might find that Table I has been cross
% referenced as Table IV-B3. 

% {\LaTeX} does not have precognitive abilities. If you put a
% \verb|\label| command before the command that updates the counter it's
% supposed to be using, the label will pick up the last counter to be
% cross referenced instead. In particular, a \verb|\label| command
% should not go before the caption of a figure or a table.

% Do not use \verb|\nonumber| inside the \verb|{array}| environment. It
% will not stop equation numbers inside \verb|{array}| (there won't be
% any anyway) and it might stop a wanted equation number in the
% surrounding equation.

% \subsection{Some Common Mistakes}\label{SCM}
% \begin{itemize}
% \item The word ``data'' is plural, not singular.
% \item The subscript for the permeability of vacuum $\mu_{0}$, and other common scientific constants, is zero with subscript formatting, not a lowercase letter ``o''.
% \item In American English, commas, semicolons, periods, question and exclamation marks are located within quotation marks only when a complete thought or name is cited, such as a title or full quotation. When quotation marks are used, instead of a bold or italic typeface, to highlight a word or phrase, punctuation should appear outside of the quotation marks. A parenthetical phrase or statement at the end of a sentence is punctuated outside of the closing parenthesis (like this). (A parenthetical sentence is punctuated within the parentheses.)
% \item A graph within a graph is an ``inset'', not an ``insert''. The word alternatively is preferred to the word ``alternately'' (unless you really mean something that alternates).
% \item Do not use the word ``essentially'' to mean ``approximately'' or ``effectively''.
% \item In your paper title, if the words ``that uses'' can accurately replace the word ``using'', capitalize the ``u''; if not, keep using lower-cased.
% \item Be aware of the different meanings of the homophones ``affect'' and ``effect'', ``complement'' and ``compliment'', ``discreet'' and ``discrete'', ``principal'' and ``principle''.
% \item Do not confuse ``imply'' and ``infer''.
% \item The prefix ``non'' is not a word; it should be joined to the word it modifies, usually without a hyphen.
% \item There is no period after the ``et'' in the Latin abbreviation ``et al.''.
% \item The abbreviation ``i.e.'' means ``that is'', and the abbreviation ``e.g.'' means ``for example''.
% \end{itemize}
% An excellent style manual for science writers is \cite{b7}.

% \subsection{Authors and Affiliations}
% \textbf{The class file is designed for, but not limited to, six authors.} A 
% minimum of one author is required for all conference articles. Author names 
% should be listed starting from left to right and then moving down to the 
% next line. This is the author sequence that will be used in future citations 
% and by indexing services. Names should not be listed in columns nor group by 
% affiliation. Please keep your affiliations as succinct as possible (for 
% example, do not differentiate among departments of the same organization).

% \subsection{Identify the Headings}
% Headings, or heads, are organizational devices that guide the reader through 
% your paper. There are two types: component heads and text heads.

% Component heads identify the different components of your paper and are not 
% topically subordinate to each other. Examples include Acknowledgments and 
% References and, for these, the correct style to use is ``Heading 5''. Use 
% ``figure caption'' for your Figure captions, and ``table head'' for your 
% table title. Run-in heads, such as ``Abstract'', will require you to apply a 
% style (in this case, italic) in addition to the style provided by the drop 
% down menu to differentiate the head from the text.

% Text heads organize the topics on a relational, hierarchical basis. For 
% example, the paper title is the primary text head because all subsequent 
% material relates and elaborates on this one topic. If there are two or more 
% sub-topics, the next level head (uppercase Roman numerals) should be used 
% and, conversely, if there are not at least two sub-topics, then no subheads 
% should be introduced.

% \subsection{Figures and Tables}
% \paragraph{Positioning Figures and Tables} Place figures and tables at the top and 
% bottom of columns. Avoid placing them in the middle of columns. Large 
% figures and tables may span across both columns. Figure captions should be 
% below the figures; table heads should appear above the tables. Insert 
% figures and tables after they are cited in the text. Use the abbreviation 
% ``Fig.~\ref{fig}'', even at the beginning of a sentence.

% \begin{table}[htbp]
% \caption{Table Type Styles}
% \begin{center}
% \begin{tabular}{|c|c|c|c|}
% \hline
% \textbf{Table}&\multicolumn{3}{|c|}{\textbf{Table Column Head}} \\
% \cline{2-4} 
% \textbf{Head} & \textbf{\textit{Table column subhead}}& \textbf{\textit{Subhead}}& \textbf{\textit{Subhead}} \\
% \hline
% copy& More table copy$^{\mathrm{a}}$& &  \\
% \hline
% \multicolumn{4}{l}{$^{\mathrm{a}}$Sample of a Table footnote.}
% \end{tabular}
% \label{tab1}
% \end{center}
% \end{table}

% \begin{figure}[htbp]
% \centerline{\includegraphics{fig1.png}}
% \caption{Example of a figure caption.}
% \label{fig}
% \end{figure}

% Figure Labels: Use 8 point Times New Roman for Figure labels. Use words 
% rather than symbols or abbreviations when writing Figure axis labels to 
% avoid confusing the reader. As an example, write the quantity 
% ``Magnetization'', or ``Magnetization, M'', not just ``M''. If including 
% units in the label, present them within parentheses. Do not label axes only 
% with units. In the example, write ``Magnetization (A/m)'' or ``Magnetization 
% \{A[m(1)]\}'', not just ``A/m''. Do not label axes with a ratio of 
% quantities and units. For example, write ``Temperature (K)'', not 
% ``Temperature/K''.

% \section*{Acknowledgment}

% The preferred spelling of the word ``acknowledgment'' in America is without 
% an ``e'' after the ``g''. Avoid the stilted expression ``one of us (R. B. 
% G.) thanks $\ldots$''. Instead, try ``R. B. G. thanks$\ldots$''. Put sponsor 
% acknowledgments in the unnumbered footnote on the first page.

% \section*{References}

% Please number citations consecutively within brackets \cite{b1}. The 
% sentence punctuation follows the bracket \cite{b2}. Refer simply to the reference 
% number, as in \cite{b3}---do not use ``Ref. \cite{b3}'' or ``reference \cite{b3}'' except at 
% the beginning of a sentence: ``Reference \cite{b3} was the first $\ldots$''

% Number footnotes separately in superscripts. Place the actual footnote at 
% the bottom of the column in which it was cited. Do not put footnotes in the 
% abstract or reference list. Use letters for table footnotes.

% Unless there are six authors or more give all authors' names; do not use 
% ``et al.''. Papers that have not been published, even if they have been 
% submitted for publication, should be cited as ``unpublished'' \cite{b4}. Papers 
% that have been accepted for publication should be cited as ``in press'' \cite{b5}. 
% Capitalize only the first word in a paper title, except for proper nouns and 
% element symbols.

% For papers published in translation journals, please give the English 
% citation first, followed by the original foreign-language citation \cite{b6}.

% \begin{thebibliography}{00}
% \bibitem{b1} G. Eason, B. Noble, and I. N. Sneddon, ``On certain integrals of Lipschitz-Hankel type involving products of Bessel functions,'' Phil. Trans. Roy. Soc. London, vol. A247, pp. 529--551, April 1955.
% \bibitem{b2} J. Clerk Maxwell, A Treatise on Electricity and Magnetism, 3rd ed., vol. 2. Oxford: Clarendon, 1892, pp.68--73.
% \bibitem{b3} I. S. Jacobs and C. P. Bean, ``Fine particles, thin films and exchange anisotropy,'' in Magnetism, vol. III, G. T. Rado and H. Suhl, Eds. New York: Academic, 1963, pp. 271--350.
% \bibitem{b4} K. Elissa, ``Title of paper if known,'' unpublished.
% \bibitem{b5} R. Nicole, ``Title of paper with only first word capitalized,'' J. Name Stand. Abbrev., in press.
% \bibitem{b6} Y. Yorozu, M. Hirano, K. Oka, and Y. Tagawa, ``Electron spectroscopy studies on magneto-optical media and plastic substrate interface,'' IEEE Transl. J. Magn. Japan, vol. 2, pp. 740--741, August 1987 [Digests 9th Annual Conf. Magnetics Japan, p. 301, 1982].
% \bibitem{b7} M. Young, The Technical Writer's Handbook. Mill Valley, CA: University Science, 1989.
% \end{thebibliography}
% \vspace{12pt}
% \color{red}
% IEEE conference templates contain guidance text for composing and formatting conference papers. Please ensure that all template text is removed from your conference paper prior to submission to the conference. Failure to remove the template text from your paper may result in your paper not being published.

\end{document}
